\documentclass[11pt,a4paper]{article}
\usepackage[margin=24mm]{geometry}
\usepackage[T1]{fontenc}
\usepackage[utf8]{inputenc}
\usepackage{lmodern}
\usepackage{amsmath,amssymb,amsthm,mathtools}
\usepackage{array,longtable,booktabs}
\usepackage{xurl}
\usepackage[unicode,hidelinks]{hyperref}
\usepackage{microtype}
\setlength{\parindent}{0pt}
\setlength{\parskip}{5pt}
\setlength{\emergencystretch}{3em}
\allowdisplaybreaks
\newcommand{\R}{\mathbb R}
\newcommand{\C}{\mathbb C}
\newcommand{\N}{\mathbb N}
\newcommand{\norm}[1]{\left\lVert#1\right\rVert}
\newcommand{\abs}[1]{\left|#1\right|}
\DeclareMathOperator{\tr}{tr}
\DeclareMathOperator{\Tr}{Tr}
\DeclareMathOperator{\diag}{diag}
\DeclareMathOperator{\rank}{rank}
\DeclareRobustCommand{\code}[1]{{\footnotesize\urlstyle{tt}\nolinkurl{#1}}}
\newtheorem{theorem}{Theorem}[section]
\newtheorem{lemma}[theorem]{Lemma}
\title{Random Matrices, Probability, and Concentration\\
Proof Routes for the Remaining Open Items}
\author{Mathematical derivations and audit for NLAlib}
\date{16 October 2026}
\begin{document}
\maketitle

\textbf{Scope.} The
\href{https://research.chen.pw/NLAlib/}{live catalogue}
on 16 October 2026 lists 15 open cards in Random matrices and four in
Probability and concentration. This document covers all 19, with the
current statements, constants, and side conditions. It is a mathematical
handoff for formalization, not a collection of newly implemented Lean proofs.
No Lean source or catalogue status was changed.

\textbf{Principal new route.} The partial-transpose Gaussian moment
bound can be proved using a finite Hermite space of dimension
$\binom{s+p}{p}$. Together with the now-formalized coefficient-block
inequalities, it proves the covariance-sensitive norm target and the
strong matrix-chaos target without their former free-probability
dependency. The latter even admits the better logarithmic exponent
$q/2$ in the displayed explicit bound. A finite scalar domination lemma
also avoids a general subgaussian-process comparison theorem.

\textbf{What is and is not closed here.} Twelve cards receive direct
derivations using finite identities and already available analysis.
Seven cards receive explicit reductions with a retained deep input:
three require the BMV trace-exponential positivity theorem, two require
two-sided affine free-spectrum comparison, and two require the global
stable-branch characterization of the block-Wishart law. The retained
inputs are specified in the relevant sections. Their consequences are
proved, but those inputs are not silently treated as existing Lean results
or claimed to have elementary proofs here. These are classical or recent
proved mathematical results; ``open'' refers to the catalogue's
formalization status.

\textbf{Snapshot and reuse.} The parsed site atlas agrees with the
repository atlas at commit
\href{https://github.com/tchen-research/NLAlib/tree/faef0a270d762f0bf0e90603544d50f289873d0e}
{\texttt{faef0a270d762f0bf0e90603544d50f289873d0e}}.
All named NLAlib reuse pointers were checked in the corresponding public
declaration index with \code{sorry=false}. Mathlib names are accompanied
by the source module inspected. This is a statement-availability check,
not a claim that proposed new applications were compiled.

\textbf{Conventions.} The operator norm is denoted by $\norm{\cdot}$ unless
another norm is specified. The adjoint is $*$; real transpose is $T$ or
$\mathsf T$. Ordinary trace is unnormalized. A normalized trace is always
written $\tr_d=d^{-1}\operatorname{Tr}$ or displayed with its factor.
Hermitian matrices form a real vector space even when their entries are
complex; covariance matching refers to every real-coordinate second
moment. All logarithms are natural.

\textbf{One law-transport principle for the whole document.} If finitely
many random objects $X_i$ on an arbitrary probability space are independent
with laws $\mu_i$, their joint law is $\bigotimes_i\mu_i$. Every measurable
matrix expression in them therefore has the same law as the corresponding
expression on that finite product. Independent copies are constructed by
adding product factors. Nonnegative integrals transfer by the map-integral
identity, and ordinary expectations transfer after integrability is checked.
All state spaces here are finite-dimensional Borel spaces; no
standard-Borel assumption on the original sample space is required.
Conditioning in the proofs can usually be implemented as iterated product
integration with deterministic matrix coefficients.

\tableofcontents
\clearpage
\section*{Coverage and dependency ledger}
\addcontentsline{toc}{section}{Coverage and dependency ledger}
``Direct'' means that the mathematical derivation is supplied here.
``BMV'', ``Free'', and ``Branch'' identify a retained input, as described
above. None of these labels changes the atlas's proof status.
\small
\begin{longtable}{@{}r p{0.64\textwidth} l l@{}}
\toprule
No. & Atlas id & Section & Route \\
\midrule
\endfirsthead
\toprule
No. & Atlas id & Section & Route \\
\midrule
\endhead
1 & \code{chi-square-neg-moment} & \ref{route:chi-square-neg-moment} & Direct \\
2 & \code{gaussian-width-def} & \ref{route:gaussian-width-def} & Direct \\
3 & \code{gaussian-width-restricted-minimum} & \ref{route:gaussian-width-restricted-minimum} & Direct \\
4 & \code{matrix-chaos-flattenings-def} & \ref{route:matrix-chaos-flattenings-def} & Direct \\
5 & \code{matrix-chaos-decoupling-upper} & \ref{route:matrix-chaos-decoupling-upper} & Direct \\
6 & \code{matrix-chaos-iterated-khintchine} & \ref{route:matrix-chaos-iterated-khintchine} & Direct \\
7 & \code{gaussian-series-covariance-sensitive-norm} & \ref{route:gaussian-series-covariance-sensitive-norm} & Direct \\
8 & \code{gaussian-trace-mgf-comparison} & \ref{route:gaussian-trace-mgf-comparison} & BMV \\
9 & \code{gaussian-proxy-eigenvalue-comparison} & \ref{route:gaussian-proxy-eigenvalue-comparison} & BMV \\
10 & \code{exchangeable-matrix-covariance} & \ref{route:exchangeable-matrix-covariance} & Direct \\
11 & \code{cauchy-transform-gaussian-universality} & \ref{route:cauchy-transform-gaussian-universality} & Direct \\
12 & \code{matrix-chaos-strong-khintchine} & \ref{route:matrix-chaos-strong-khintchine} & Direct \\
13 & \code{sample-covariance-fourth-moment-lower} & \ref{route:sample-covariance-fourth-moment-lower} & BMV \\
14 & \code{gaussian-free-spectrum-two-sided} & \ref{route:gaussian-free-spectrum-two-sided} & Free \\
15 & \code{gaussian-quadratic-free-spectrum} & \ref{route:gaussian-quadratic-free-spectrum} & Free \\
16 & \code{block-wishart-variational-edges} & \ref{route:block-wishart-variational-edges} & Branch \\
17 & \code{block-wishart-log-potential} & \ref{route:block-wishart-log-potential} & Branch \\
18 & \code{gaussian-series-partial-transpose-norm} & \ref{route:gaussian-series-partial-transpose-norm} & Direct \\
19 & \code{gaussian-cayley-fourier-block-law} & \ref{route:gaussian-cayley-fourier-block-law} & Direct \\
\bottomrule
\end{longtable}
\normalsize
\clearpage
\section{Four targets needing only finite comparisons and linear algebra}

\subsection{Negative chi-square moments: reuse the endpoint and interpolate}
\label{route:chi-square-neg-moment}
\textbf{Atlas id:} \code{chi-square-neg-moment}.
Let $S=\sum_{i=1}^d g_i^2$, with independent standard real Gaussian
coordinates and $d\ge1$. The existing density or Gamma-integral calculation
gives, for $0\le q<d/2$,
\[
 \mathbb E S^{-q}
 =\frac1{2^{d/2}\Gamma(d/2)}
      \int_0^\infty s^{d/2-q-1}e^{-s/2}\,ds
 =\frac{\Gamma(d/2-q)}{2^q\Gamma(d/2)}.                 \tag{W1}
\]
The exponent at zero is greater than $-1$, which is the integrability
condition; the exponential controls infinity. Also $S>0$ almost surely,
so the value assigned to an inverse on the null event $S=0$ is irrelevant.
For $q=0$ the answer is $1$. Two uses of the Gamma recurrence at $q=2$
give $\mathbb E S^{-2}=1/((d-2)(d-4))$ when $d\ge5$.
These statements are already proved in the current declaration index.

The remaining uniform bound requires no further Gamma estimate. Put
$r=(d-1)/2$ and suppose $0<q\le r$. Concavity of $u\mapsto u^{q/r}$ on
the nonnegative half-line, on a probability space, gives
\[
 \mathbb E S^{-q}\le(\mathbb E S^{-r})^{q/r},\qquad
 (\mathbb E S^{-q})^{1/q}\le(\mathbb E S^{-r})^{1/r}.    \tag{W2}
\]
For $q=r$ this is equality. Integrability of the smaller power also follows
from $u^{q/r}\le1+u$, so Jensen has all its hypotheses.
For $d\ge5$, the already proved endpoint inequality
$\mathbb E S^{-r}<(3/d)^r$ finishes the desired bound $3/d$.
For the only remaining admissible dimensions, (W1) evaluates the endpoint
exactly:
\[
\begin{array}{c|c|c}
 d&r& d(\mathbb ES^{-r})^{1/r}\\\hline
 2&1/2&\pi\\
 3&1&3\\
 4&3/2&2\pi^{1/3}.
\end{array}
\]
Each entry is below $4$ (use $\pi<4$, and $\pi<8$ in the last case).
For $d=1$, the hypotheses $q>0$ and $d\ge2q+1$ are incompatible.
This proves the complete card, including the distinction between mere
integrability $2q<d$ and the stronger uniform-bound condition $2q+1\le d$.
The latter cannot be dropped: the Gamma factor in (W1) diverges as
$q\uparrow d/2$ at fixed $d$.

\paragraph{Verified reuse.}
The exact formula, integrability, and large-dimension endpoint are
\code{NLAlib.integral_rpow_inv_sum_sq_gaussianVector},
\code{NLAlib.integrable_rpow_inv_sum_sq_gaussianVector}, and
\code{NLAlib.integral_rpow_inv_sum_sq_gaussianVector_lt}.
Even the interpolation step (W2) already exists as
\code{NLAlib.integral_norm_rpow_inv_le}, applied to the nonnegative random
variable $S^{-1}$. The second inverse moment is
\code{NLAlib.integral_inv_sum_sq_pow_two_gaussianVector}.
Thus the missing formal work for this card is principally the endpoint
assembly, dimensions $2,3,4$, and transport along the Gaussian law.

\subsection{Gaussian width: define the support function first}
\label{route:gaussian-width-def}
\textbf{Atlas id:} \code{gaussian-width-def}.
For a nonempty compact $T\subset\R^n$, put
$R_T=\max_{x\in T}\norm x_2$ and
\[
 h_T(z)=\max_{x\in T}\langle z,x\rangle,\qquad
 w(T)=\mathbb E h_T(g),\quad g\sim N(0,I_n).
\]
Compactness supplies both maxima. The pointwise inequalities
\[
 |h_T(z)|\le R_T\norm z_2,\qquad
 |h_T(z)-h_T(z')|\le R_T\norm{z-z'}_2                 \tag{W3}
\]
follow by bounding each linear functional and taking suprema. In particular,
the support function is continuous and integrable, since
$\mathbb E\norm g_2\le\sqrt{\mathbb E\norm g_2^2}=\sqrt n$.
This avoids any measurable-maximizer theorem or selection of a random
maximizing point. If $x_0\in T$, then $h_T(g)\ge\langle g,x_0\rangle$,
so $0\le w(T)\le R_T\sqrt n$.

The basic identities now hold before integration:
\[
 h_{T+a}(z)=h_T(z)+\langle z,a\rangle,\quad
 h_{UT}(z)=h_T(U^Tz),\quad
 h_{\operatorname{conv}T}(z)=h_T(z)
                     =h_{\overline{\operatorname{conv}}T}(z).    \tag{W4}
\]
The last equality follows first for each finite convex combination, then
by continuity for its closure. It proves the corresponding expected
identities without an additional compactness argument: the same support
function is finite and integrable. Alternatively the closed convex hull
is closed and bounded, hence compact in finite dimension.
Centering and rotational invariance of $g$ give
\[
 w(T+a)=w(T),\qquad w(UT)=w(T)\quad(U^TU=I),\qquad
 w(\operatorname{conv}T)=w(T).
\]
Also $w(cT)=c\,w(T)$ for $c\ge0$, and inclusion of sets increases width.
Width here has no absolute value inside the supremum. In particular every
singleton has width zero, regardless of its distance from the origin.

If $L\ne\{0\}$ is a subspace and $T$ is its unit sphere, then
$h_T(z)=\norm{P_Lz}_2$: Cauchy--Schwarz gives the upper bound, and normalize
$P_Lz$ to attain it when nonzero. When it is zero every unit vector in $L$
attains zero. Choosing an orthonormal basis of $L$ shows
$w(T)=\mathbb E\norm{g_{\dim L}}_2$. The unit sphere of the zero subspace
is empty and is intentionally outside this definition. If $n=0$, the only
nonempty compact subset is $\{0\}$ and its width is zero.

\paragraph{Verified reuse.}
Use \code{NLAlib.integrable_of_lipschitz_gaussianVector} with (W3),
\code{NLAlib.gaussianVector_map_toLp_eq_stdGaussian}, and
\code{NLAlib.hasLaw_gaussianMatrix_orthogonal} with one matrix column.
The compact support-function construction is the small deterministic
interface to add; a new theory of Gaussian processes is unnecessary.

\subsection{Restricted minimum: finite Gordon arrays and deterministic nets}
\label{route:gaussian-width-restricted-minimum}
\textbf{Atlas id:} \code{gaussian-width-restricted-minimum}.
Let $m,n\ge1$, let $G$ be a standard real $m\times n$ Gaussian matrix,
and let $\varnothing\ne T\subset S^{n-1}$ be compact. Define
$f(G)=\min_{x\in T}\norm{Gx}_2$. For every $A,B$,
\[
 |f(A)-f(B)|\le\norm{A-B}_2\le\norm{A-B}_F.             \tag{W5}
\]
One proves this by testing the minimizer of one matrix in the other and
then reversing the roles. Thus $f$ is measurable and integrable, without
choosing those minimizers measurably.

Here is the full reduction to the already proved \emph{finite} Gordon
theorem. Choose finite nonempty sets $T_0\subset T$, $V_0\subset S^{m-1}$.
Let $a\sim N(0,I_m)$, $b\sim N(0,I_n)$ and $\gamma\sim N(0,1)$ be
independent of each other and of $G$. For $x\in T_0$, $y\in V_0$, set
\[
 U_{xy}=\langle a,y\rangle+\langle b,x\rangle,\qquad
 V_{xy}=\langle Gx,y\rangle+\gamma.
\]
Both arrays are centered jointly Gaussian. For fixed $x$, their increment
variances in $y$ agree. For $x,x'\in T_0$, $y,y'\in V_0$, putting
$u=\langle x,x'\rangle$, $v=\langle y,y'\rangle$, gives
\[
 \mathbb E(U_{xy}-U_{x'y'})^2=4-2u-2v,\quad
 \mathbb E(V_{xy}-V_{x'y'})^2=2-2uv,
\]
whose difference is $2(1-u)(1-v)\ge0$. The finite Gordon inequality in
the indexed orientation therefore yields
\[
 \mathbb E\max_{y\in V_0}\langle a,y\rangle-w(T_0)
 \le\mathbb E\min_{x\in T_0}\max_{y\in V_0}\langle Gx,y\rangle. \tag{W6}
\]
Here $\mathbb E\min_x\langle b,x\rangle=-w(T_0)$ follows from symmetry
$b\overset{\rm law}=-b$, and the common scalar $\gamma$ integrates to zero.

For $\epsilon>0$ choose deterministic finite $\epsilon$-nets
$T_\epsilon\subset T$ and $V_\epsilon\subset S^{m-1}$. The support-function
errors are at most $\epsilon\norm b$ and $\epsilon\norm a$. Moreover
\[
 \left|\min_{x\in T_\epsilon}\max_{y\in V_\epsilon}\langle Gx,y\rangle
                  -f(G)\right|\le2\epsilon\norm G_F.            \tag{W7}
\]
For the inner maximum approximate a unit vector attaining $\norm{Gx}$;
for the outer minimum approximate a minimizing point of $T$. These are
deterministic inequalities for each matrix, not random choices used as
measurable functions. Since all three norm expectations are finite,
(W6)--(W7) pass to the limit simply by bounding expectation differences.
They prove
\[
 \mathbb Ef(G)\ge\mathbb E\norm a_2-w(T).                       \tag{W8}
\]

The already proved Gordon singular-value theorem with one column gives
$\mathbb E\norm a_2\ge\sqrt m-1$. An elementary independent check is:
for $m\ge2$, the chi-density and Gamma recurrence give
$\mathbb E\norm a=(m-1)\mathbb E\norm a^{-1}$, while Cauchy--Schwarz
gives $\mathbb E\norm a\,\mathbb E\norm a^{-1}\ge1$; hence
$\mathbb E\norm a\ge\sqrt{m-1}\ge\sqrt m-1$. For $m=1$, nonnegativity
is sufficient for the asserted weaker bound.
Gaussian concentration applied to $-f$, whose Lipschitz constant is one
by (W5), now gives for $t\ge0$
\[
 \boxed{\mathbb P\{f(G)<\sqrt m-1-w(T)-t\}\le e^{-t^2/2}.}       \tag{W9}
\]
The event is contained in $\{f(G)\le\mathbb Ef(G)-t\}$.
If the displayed threshold is nonpositive, it is empty; the strict event
avoids making any assertion about an atom at zero.

\paragraph{Verified reuse.}
The exact finite array interface is \code{NLAlib.gordon_minimax_inequality}.
The required tail is \code{NLAlib.gaussian_concentration}, or its
arbitrary-space transport \code{NLAlib.gaussian_concentration_of_hasLaw}.
Use \code{NLAlib.gordon_extreme_singular_values} with one column for the
Gaussian radius bound. Only compact nets and (W7) bridge the existing finite
comparison to the set-indexed target; no separable-process extension is needed.

\subsection{Gaussian Cayley matrices: explicit finite Fourier certificates}
\label{route:gaussian-cayley-fourier-block-law}
\textbf{Atlas id:} \code{gaussian-cayley-fourier-block-law}.
Let $G$ be a finite group of order $N\ge1$, and let $\pi$ run through a
complete family of inequivalent irreducible unitary complex representations,
of dimensions $d_\pi$. The route below proves the required finite Fourier
decomposition directly; it does not need a general Peter--Weyl theorem.

\paragraph{Finite Schur averaging.}
For a map $M:\C^{d_\tau}\to\C^{d_\pi}$, put
\[
 \mathcal A_{\pi\tau}(M)=\frac1N\sum_{g\in G}\pi(g)M\tau(g)^*.
\]
Changing the summation variable shows this is an intertwiner. Its kernel
and range are invariant. Thus a nonzero intertwiner between irreducibles
is an isomorphism, so it vanishes when $\pi\ne\tau$. For $\pi=\tau$,
choose an eigenvalue of the intertwiner over $\C$; the nonzero kernel of
$\mathcal A-\lambda I$ is invariant and therefore is the whole space.
The intertwiner is scalar. Taking traces identifies the scalar:
$\mathcal A_{\pi\pi}(M)=\tr(M)I/d_\pi$.
Applying this to a matrix unit gives the coefficient identity
\[
 \frac1N\sum_{g\in G}\pi(g)_{ia}\overline{\tau(g)_{jb}}
 =\begin{cases}d_\pi^{-1}\delta_{ij}\delta_{ab},&\pi=\tau,\\
                 0,&\pi\ne\tau.\end{cases}                    \tag{W10}
\]
There is no probability in this step.

\paragraph{Completeness without a dimension formula as an extra axiom.}
On $\C^G$, with counting inner product, define
$\psi_{\pi ab}(g)=\sqrt{d_\pi/N}\,\pi(g^{-1})_{ab}$.
Equation (W10), with $g\mapsto g^{-1}$, makes these functions orthonormal.
Their span $M$ is invariant under $\rho(h)f(g)=f(h^{-1}g)$, since
\[
 \rho(h)\psi_{\pi ab}=\sum_c\psi_{\pi ac}\pi(h)_{cb}.           \tag{W11}
\]
Their orthogonal complement is invariant as well, because $\rho(h)$ is
unitary and its inverse is another group element.
If $M^\perp\ne0$, choose a nonzero invariant subspace $W\subseteq M^\perp$
of minimal dimension. Its restricted representation is irreducible, hence
equivalent to some $\pi$ by completeness of the given list. Let
$J:\C^{d_\pi}\to W$ be an invertible intertwiner. Evaluation at the
identity $e\in G$ gives, for $v\in\C^{d_\pi}$,
\[
 (Jv)(g)=(\rho(g^{-1})Jv)(e)=(J\pi(g^{-1})v)(e).
\]
The last expression is a linear combination of the coefficients
$\pi(g^{-1})_{ab}$, so $Jv\in M$, a contradiction. Thus the orthonormal
family is a basis, and $\sum_\pi d_\pi^2=N$ follows from its size.
With columns $\psi_{\pi ab}$, the unitary matrix $U$ satisfies, by (W11),
\[
 U^*\rho(h)U=\bigoplus_\pi I_{d_\pi}\otimes\pi(h).              \tag{W12}
\]
The index $a$ is the multiplicity coordinate; $b$ is the representation
coordinate. Stating this order explicitly avoids a tensor-transpose error.

\paragraph{Gaussian independence by realification of one unitary.}
For $z_g=\xi_g+i\eta_g$, all $2N$ real coordinates independent $N(0,1)$,
define
\[
 B_\pi=\sum_gz_g\pi(g),\qquad Z_\pi=\sqrt{d_\pi/N}\,B_\pi.
\]
The transformation from $(z_g)$ to the concatenated entries of $(Z_\pi)$
has coefficient rows $\sqrt{d_\pi/N}\,\pi(g)_{ab}$. They are orthonormal
by (W10), and their number is $N$, so this transformation is unitary.
If its complex matrix is $F=P+iQ$, the real transformation of real and
imaginary coordinates is
\[
 \mathcal R(F)=\begin{bmatrix}P&-Q\\Q&P\end{bmatrix},\qquad
 \mathcal R(F)^T\mathcal R(F)=I.
\]
Real Gaussian rotational invariance therefore says that all the resulting
real and imaginary coordinates are again independent standard normals.
This proves independence of every entry across every $Z_\pi$, without
developing a new complex Gaussian covariance API.

Finally (W12) holds pointwise for every coefficient vector, so
\[
 Z_G=\sum_gz_g\rho(g)
 =U\left[\bigoplus_\pi I_{d_\pi}\otimes
                   \sqrt{N/d_\pi}\,Z_\pi\right]U^*.
\]
This is stronger than the requested identity in law, with the blocks
defined on the original space. Orthogonal direct sums and unitary invariance
give the claimed norm law
$\norm{Z_G}_2\overset{\rm law}=\max_\pi\sqrt{N/d_\pi}\norm{Z_\pi}_2$.
The convention is $\mathbb E|z_g|^2=2$. Real-only Gaussian coefficients do
not satisfy the stated independence conclusion; conjugate representations
then impose relations. For the trivial group, $N=d_\pi=1$, and every
identity reduces to the single circular Gaussian variable.

\paragraph{Verified reuse and what remains to package.}
Mathlib already supplies
\code{Representation.IsIrreducible.bijective_or_eq_zero} and
\code{Representation.IsIrreducible.algebraMap_intertwiningMap_bijective_of_isAlgClosed}
in \code{Mathlib.RepresentationTheory.Irreducible}.
Their local source was inspected. These replace the two Schur-lemma
arguments. Equation (W10), the finite matrix-coefficient basis (W11),
and the realification calculation are the useful small bridges to add.
The probability input is the already proved
\code{NLAlib.hasLaw_gaussianMatrix_orthogonal} on a real $2N$-dimensional
one-column Gaussian matrix. Thus neither infinite-dimensional harmonic
analysis nor an independent complex normal-distribution library is needed.


\section{Matrix chaos through finite coefficient arrays}

This section covers five open cards.  The principal change of route is that
the new partial-transpose estimate can be proved entirely on a finite
Hermite space.  It supplies the strong linear estimate needed for chaos,
so that the strong chaos theorem need not wait for a free-probability
comparison theorem.  All laws may be realized on an arbitrary probability
space: independence identifies their finite-dimensional joint laws with
product laws, and the calculations below take place on those product laws.

\subsection{Coefficient flattenings are reindexed finite matrices}
\label{route:matrix-chaos-flattenings-def}
\noindent\code{matrix-chaos-flattenings-def}.

Let $q,m,d_1,d_2\geq1$, and let
$A_{i_1,\ldots,i_q}\in\R^{d_1\times d_2}$.  Regard the coefficients as the
array
\[
 a(i_1,\ldots,i_q,a,b)=(A_{i_1,\ldots,i_q})_{ab}.
\]
For $S\subseteq[q]$, define the ordinary finite matrix $F_S$ by
\[
 (F_S)_{((i_t)_{t\in S},a),((i_t)_{t\notin S},b)}
       =a(i_1,\ldots,i_q,a,b),\qquad
 \sigma(A)=\max_{S\subseteq[q]}\norm{F_S}.
\]
Products over an empty set are singleton index types.  Thus $F_S$ has
$d_1m^{|S|}$ rows and $d_2m^{q-|S|}$ columns.  No tensor-product
construction is needed: an entry is obtained by the unique assembly of
its row and column coordinates.  Changing the ordering within either
index set multiplies by permutation matrices, so preserves the norm.

For independent coordinate groups $h^{(1)},\ldots,h^{(q)}$, define
\[
 Y=\sum_{i_1,\ldots,i_q=1}^m
       \left(\prod_{t=1}^q h^{(t)}_{i_t}\right)A_{i_1,\ldots,i_q}.
\]
Coincident values of indices in different groups are allowed.  For $q=1$,
direct multiplication of horizontal and vertical block matrices gives
\[
 \norm{F_\varnothing}^2=\left\|\sum_iA_iA_i^*\right\|,
 \qquad
 \norm{F_{\{1\}}}^2=\left\|\sum_iA_i^*A_i\right\|.
\]
The covariance flattenings needed later are
\[
 (V_S)_{(i_t)_{t\in S},((i_t)_{t\notin S},a,b)}
       =a(i_1,\ldots,i_q,a,b),\qquad
 v(A)=\max_{\varnothing\ne S\subseteq[q]}\norm{V_S}.
\]
The row subset is nonempty, and both original matrix coordinates are on
the column side.  These are different from flattenings that repeat a
coordinate on both sides; no such overlapping flattening is used here.

\paragraph{Finite bookkeeping lemma.}
Freeze the last chaos coordinate $j=i_q$ and write
$Y=\sum_jh_j^{(q)}B_j$, where $B_j$ is a chaos of degree $q-1$.
Form its horizontal, vertical, and vectorization matrices
\[
 R=[B_1\ \cdots\ B_m],\qquad
 C=[B_1^*\ \cdots\ B_m^*]^*,\qquad
 V_{j,(a,b)}=(B_j)_{ab}.
\]
Each is a lower-degree chaos obtained by moving $j$ into its matrix
indices.  Its coefficient parameters satisfy
\begin{align*}
 &\sigma(R),\sigma(C)\leq\sigma(A),\qquad
 v(R),v(C)\leq v(A),\\
 &\sigma(V)\leq v(A).
\end{align*}
Here a symbol such as $\sigma(R)$ means the parameter of the coefficient
array of that lower-degree chaos.  To verify the first line, append $j$
to the column or row subset of an $F$-flattening; for a $v$-flattening,
the new matrix coordinates, including $j$, all go to its column side.
For the second line, every variance flattening of $V$ has row set
$\{q\}\cup S$, which is nonempty, and has $a,b$ on its column side.
These are entrywise identifications up to permutations, not norm estimates
for an abstract tensor product.  For degree zero set $\sigma(B)=\norm B$
and $v(B)=0$ only for purposes of induction.

\paragraph{Sources and reuse.}
The definitions agree with
\href{https://arxiv.org/html/2412.18468v2}{Bandeira--Lucca--Nizi\'c-Nikolac--van Handel,
Sections 2.1--2.2}.  Existing block multiplication, permutation isometries,
and \code{NLAlib.dilation_identities} suffice for the finite matrix
interfaces; a new general-purpose tensor library is unnecessary.

\subsection{One-sided decoupling by a finite coloring average}
\label{route:matrix-chaos-decoupling-upper}
\noindent\code{matrix-chaos-decoupling-upper}.

\begin{theorem}
Suppose $h_1,\ldots,h_m$ are independent centered integrable real variables,
and $A_{i_1,\ldots,i_q}=0$ whenever two indices coincide.  If $X$ is the
coupled chaos and $Y$ its fully decoupled version, then
\[
                     \mathbb E\norm X\leq q^q\mathbb E\norm Y.
\]
No symmetry under permutations of the $q$ coefficient indices is required.
\end{theorem}
\begin{proof}
Average over the finite set of colorings $c:[m]\to[q]$, giving each
coloring mass $q^{-m}$.  Put
\[
 X_c=\sum_{i_1,\ldots,i_q}
       \prod_{t=1}^q\bigl(1_{\{c(i_t)=t\}}h_{i_t}\bigr)
       A_{i_1,\ldots,i_q}.
\]
For a nonzero coefficient all indices are distinct, so exactly the
fraction $q^{-q}$ of colorings retains its term.  Hence
$\mathbb E_c X_c=q^{-q}X$, and the triangle inequality for a finite
convex combination gives
\[
 \mathbb E\norm X\leq q^q\mathbb E_c\mathbb E\norm{X_c}.
\]
Fix $c$.  In the product realization of $Y$, retain the coordinates
$h_i^{(t)}$ with $c(i)=t$ and integrate every other coordinate.  Centering
makes every term containing a nonretained coordinate vanish.  The
resulting function $Y_c$ has exactly the same law as $X_c$: the retained
coordinates come from disjoint independent classes.  Pointwise in the
retained coordinates, the norm of the integrated matrix is at most the
integral of its norm.  Thus $\mathbb E\norm{X_c}\leq\mathbb E\norm Y$.
This proves the assertion using only finite-product integration and
finite sums; a general conditional-expectation interface is optional.

All terms are integrable.  In $X$ the nonzero terms use distinct
independent factors, and in $Y$ they use independent groups, so the
expectation of each absolute scalar product is a product of finite first
moments.  If $q>m$, $X=0$ because no distinct tuple exists.  For $q=1$
the assertion is equality.  The zero coefficient family is immediate.
\end{proof}

This is the coloring proof of the upper direction of
\href{https://arxiv.org/html/2412.18468v2}{the matrix chaos decoupling theorem}.
It gives the numerical constant $q^q$ directly.  It proves no reverse
inequality: antisymmetric coefficients can make $X$ identically zero
while $Y$ is nonzero.

\subsection{Partial transpose: a finite Hermite matrix proof}
\label{route:gaussian-series-partial-transpose-norm}
\noindent\code{gaussian-series-partial-transpose-norm}.

For Hermitian $A_1,\ldots,A_s\in\C^{d\times d}$ with $s,d\geq1$, put
\[
 X=\sum_{i=1}^sg_iA_i,\quad
 S=\sum_iA_i^2,\quad \sigma^2=\norm S,\quad
 T=[A_jA_i]_{i,j=1}^s,\quad \kappa^2=\lambda_{\max}(T).
\]
The matrix $T$ is Hermitian, but is generally not positive semidefinite.
The already established coefficient-block theorem gives
\[
 0\leq\kappa\leq\sigma,
 \qquad T\preceq\kappa^2I,
 \qquad \kappa\leq v,
 \quad
 v^2=\left\|\sum_i\operatorname{vec}(A_i)
                       \operatorname{vec}(A_i)^*\right\|.
\]
The use of the largest eigenvalue, rather than the absolute largest
eigenvalue, is essential in the following estimate.

\begin{theorem}\label{thm:finite-hermite-partial-transpose}
For every integer $p\geq1$,
\[
 \left(\mathbb E\,d^{-1}\tr X^{2p}\right)^{1/(2p)}
       \leq2\sqrt{\sigma^2+(p-1)\kappa^2}.
\]
Moreover,
\[
 \mathbb E\norm X\leq2\sigma+6\kappa\sqrt{\log d}.
\]
For $d\geq2$ the slightly more precise intermediate bound is
$2\sigma+4\kappa\sqrt{\log(2d)}$.
\end{theorem}
\begin{proof}
Fix $p$ and use the finite orthonormal basis
\[
 \{e_\alpha:\alpha\in\N^s,\ |\alpha|\leq p\}
 \quad\hbox{of }H_p,
 \qquad \dim H_p=\binom{s+p}{p}.
\]
Define finite matrices
\[
 C_i e_\alpha=
 \begin{cases}
 \sqrt{\alpha_i+1}\,e_{\alpha+e_i},&|\alpha|<p,\\
 0,&|\alpha|=p,
 \end{cases}
 \qquad
 C_i^*e_\alpha=\sqrt{\alpha_i}\,e_{\alpha-e_i},
\]
with zero understood when $\alpha_i=0$.  Let $P$ project onto degrees
at most $p-1$, and let $N e_\alpha=|\alpha|e_\alpha$.  Direct basis
calculation gives
\[
 C_i^*C_jP=(\delta_{ij}I+C_jC_i^*)P,
 \qquad \sum_iC_iC_i^*=N.
\]
The first identity is deliberately restricted to $P$; the unqualified
commutation relation would be false at the cutoff.

Set $L=\sum_iA_i\otimes C_i$, and let
$Bz=((I_d\otimes C_i^*)z)_i$.  On $\C^d\otimes PH_p$ the preceding
identities give the finite factorization
\begin{align*}
 L^*L
 &=S\otimes I+\sum_{i,j}A_iA_j\otimes C_jC_i^*\\
 &=S\otimes I+B^*(T\otimes I)B
 \preceq\sigma^2I+\kappa^2(I_d\otimes N)\\
 &\preceq\bigl(\sigma^2+(p-1)\kappa^2\bigr)I.
\end{align*}
The identification of the block index in $B^*(T\otimes I)B$ follows by
interchanging $i,j$ in its sum.  Also $L$ annihilates degree $p$, and
$L^*L$ preserves degree.  Consequently the same norm bound applies on
the whole finite space:
\[
 \norm L^2\leq\sigma^2+(p-1)\kappa^2,
 \qquad
 \norm{L+L^*}\leq2\sqrt{\sigma^2+(p-1)\kappa^2}.
\]

Here is a finite verification of the probabilistic identity.  Define
the monic probabilists' Hermite polynomials by
$H_0=1$, $H_1=x$, and
$H_{k+1}=xH_k-kH_{k-1}$.  The scalar Gaussian moment recurrence implies,
by linearity on monomials,
\[
             \mathbb E[g f(g)]=\mathbb E[f'(g)]
             \quad\hbox{for every polynomial }f.
\]
This does not apply a bounded-test-function Stein lemma to an unbounded
polynomial.  It follows directly from the finite moment formulas.
Induction gives $H_k'=kH_{k-1}$ and
$\mathbb E[H_k(g)f(g)]=\mathbb E[f^{(k)}(g)]$.
Thus $h_k=H_k/\sqrt{k!}$ are orthonormal and
\[
 gh_k=\sqrt{k+1}\,h_{k+1}+\sqrt{k}\,h_{k-1}.
\]
Independence makes the finite family
$h_\alpha(g)=\prod_i h_{\alpha_i}(g_i)$ orthonormal.  Identify
$e_\alpha$ with $h_\alpha$.  Starting at the constant polynomial, the
first $p$ applications of multiplication by $X(g)$ never require a
degree exceeding $p$.  Hence, for each standard basis vector $u_a$,
\[
 (L+L^*)^p(u_a\otimes e_0)
       \quad\hbox{represents exactly}\quad X(g)^p u_a.
\]
Taking its squared norm, summing over $a$, and using that $X$ is
Hermitian proves
\[
 \mathbb E\,d^{-1}\tr X^{2p}
   =d^{-1}\sum_{a=1}^d\norm{(L+L^*)^p(u_a\otimes e_0)}^2
   \leq\norm{L+L^*}^{2p}.
\]
Only a finite orthonormal family is used.  No completeness assertion for
all Hermite polynomials, unbounded operator, or infinite Fock space is
needed.

For the norm estimate suppose first $d\geq2$, put
$\ell=\log(2d)>1$, and assume $\kappa>0$.  Choose the integer
\[
              p=\left\lceil\ell+\frac{\sigma\sqrt\ell}{\kappa}
                 \right\rceil.
\]
The scalar inequalities $p\geq\ell$ and $e^u\leq1+2u$ for
$0\leq u\leq1/2$ imply $d^{1/(2p)}\leq1+\ell/p$.
Since $\kappa\leq\sigma$,
\begin{align*}
 \sigma^2+p\kappa^2
 &\leq\sigma^2+\ell\kappa^2+\sigma\kappa\sqrt\ell+\kappa^2
 \leq(\sigma+\kappa\sqrt\ell)^2,\\
 \frac\ell p(\sigma+\kappa\sqrt\ell)&\leq\kappa\sqrt\ell.
\end{align*}
For a Hermitian matrix, $\norm X^{2p}\leq\tr X^{2p}$, and the scalar
$L^{2p}$ bound gives
\begin{align*}
 \mathbb E\norm X
 &\leq2d^{1/(2p)}\sqrt{\sigma^2+(p-1)\kappa^2}\\
 &\leq2\sigma+4\kappa\sqrt{\log(2d)}
 \leq2\sigma+6\kappa\sqrt{\log d}.
\end{align*}
The last step uses $\log(2d)\leq2\log d$ and $4\sqrt2<6$.
If $\kappa=0$, take $p\to\infty$ in the moment bound to obtain
$\mathbb E\norm X\leq2\sigma$; alternatively the diagonal blocks of
$T\preceq0$ force every $A_i=0$.  If $d=1$, $X$ is a scalar real
Gaussian of variance $\sigma^2$, and
$\mathbb E|X|=\sqrt{2/\pi}\,\sigma\leq2\sigma$.
All-zero coefficients are included in these arguments.
\end{proof}

\paragraph{Reuse and normalization audit.}
The declaration index at the recorded snapshot marks each of
\code{NLAlib.coefficientBlock_variance_covariance_bounds},
\code{NLAlib.integral_prod_gaussianVector_eq_sum_pairings}, and
\code{NLAlib.integral_abs_rpow_gaussianReal} as
\code{sorry = false}.  The first supplies the coefficient bounds; the
last two supply Gaussian moments and their arbitrary-space transport.
The already proved
\code{NLAlib.integral_re_trace_gaussianSeries_pow_rpow_le_of_hasLaw}
also supplies ordinary Gaussian noncommutative Khintchine directly.
The finite proof above strengthens it with $\kappa$.
The source is
\href{https://arxiv.org/html/2610.01982v1}{Bandeira--Kunisky--Nizi\'c-Nikolac--Pesenti--Wang,
Section 2 and Appendix A.2}.  Its trace is normalized; for the library's
unnormalized trace multiply the rooted moment bound by $d^{1/(2p)}$.

\subsection{Iterated Khintchine without a subgaussian-process theorem}
\label{route:matrix-chaos-iterated-khintchine}
\noindent\code{matrix-chaos-iterated-khintchine}.

\begin{lemma}[Scalar domination, followed by finite contraction]
\label{lem:subgaussian-gaussian-domination}
Let $h_i$ be independent centered real variables satisfying
$\mathbb E e^{h_i^2/K^2}\leq2$ for $K>0$.  For arbitrary vectors $a_i$
in a finite-dimensional normed space,
\[
          \mathbb E\left\|\sum_i h_i a_i\right\|
        \leq 10K\,\mathbb E\left\|\sum_i g_i a_i\right\|.
\]
\end{lemma}
\begin{proof}
Introduce independent copies $h_i'$.  Centering and Jensen give
$\mathbb E\norm{\sum_i h_i a_i}
 \leq\mathbb E\norm{\sum_i\xi_i a_i}$, with
$\xi_i=h_i-h_i'$ symmetric and independent.  Since
$(h-h')^2\leq2h^2+2(h')^2$, Cauchy--Schwarz and independence give
\[
 \mathbb E e^{\xi_i^2/(4K^2)}\leq2,
 \qquad \mathbb P\{|\xi_i|>t\}\leq2e^{-t^2/(4K^2)}.
\]
For a standard Gaussian $g$ and $x\geq0$,
\[
 \mathbb P\{|g|>x\}
 \geq2\int_x^{x+1}\frac{e^{-u^2/2}}{\sqrt{2\pi}}\,du
 \geq\sqrt{2/\pi}\,e^{-(x+1)^2/2}.
\]
At $t=4K(1+x)$, the preceding upper bound on the tail of $|\xi_i|$
is no larger than this lower bound.  Indeed
$2e^{-4(1+x)^2}\leq\sqrt{2/\pi}e^{-(1+x)^2/2}$.
For $t<4K$, the tail of $4K(1+|g|)$ is one.  Therefore
\[
                  |\xi_i|\ \hbox{is stochastically dominated by}
                     \ 4K(1+|g|).
\]

Here is the only contraction principle required.  For fixed $z,a$,
\[
        t\longmapsto\tfrac12(\norm{z+ta}+\norm{z-ta})
\]
is even and convex, hence nondecreasing for $t\geq0$.  Nonnegative
variables ordered by all their upper tails have ordered expectations
against every nondecreasing continuous nonnegative function.  A direct
proof approximates that function from below by increasing step
functions and applies monotone convergence.  Apply this observation
one coordinate at a time, averaging independent signs.  Since a
symmetric variable has the law of an independent sign times its
absolute value, this proves
\begin{align*}
 \mathbb E\norm{\sum_i\xi_i a_i}
 &\leq4K\,\mathbb E\norm{\sum_i\varepsilon_i(1+|g_i|)a_i}\\
 &\leq4K\left(\mathbb E\norm{\sum_i\varepsilon_i a_i}
                  +\mathbb E\norm{\sum_i g_i a_i}\right).
\end{align*}
Finally Jensen over the Gaussian magnitudes, conditional on the signs,
gives
\[
 \sqrt{2/\pi}\,\mathbb E\norm{\sum_i\varepsilon_i a_i}
       \leq\mathbb E\norm{\sum_i g_i a_i}.
\]
Thus the constant $4(1+\sqrt{\pi/2})<10$ works.  No coupling theorem,
generic chaining, or comparison of arbitrary subgaussian processes is
needed.  All needed integrability follows from the assumed exponential
square moment.
\end{proof}

For rectangular coefficients $B_j\in\R^{r\times c}$, let
\[
 \sigma_B=\max\left\{\left\|\sum_jB_jB_j^*\right\|^{1/2},
                    \left\|\sum_jB_j^*B_j\right\|^{1/2}\right\}.
\]
Hermitian dilation and either the existing Gaussian Khintchine theorem
or Theorem~\ref{thm:finite-hermite-partial-transpose} with
$\kappa\leq\sigma$ imply
\[
 \mathbb E\norm{\sum_jg_jB_j}
          \leq5\sqrt{\log(2(r+c))}\,\sigma_B.
\]
For completeness choose $p=\lceil\log(2(r+c))\rceil$.  Then
$p\leq2\log(2(r+c))$, $(r+c)^{1/(2p)}\leq\sqrt e$, and the bound
$2\sqrt p\,\sigma_B$ for the normalized rooted moment proves the
claim because $2\sqrt{2e}<5$.
Lemma~\ref{lem:subgaussian-gaussian-domination} gives the linear
subgaussian estimate with constant $50K$.

Put
\[
                  \mathcal L=\log(4d_1d_2m^q)>1.
\]
At every step of the flattening recursion the product of the two matrix
dimensions and all remaining chaos dimensions equals $d_1d_2m^q$.
Consequently every logarithm $\log(2(r+c))$ encountered is at most
$\mathcal L$.  At one step use
$\sigma_B\leq\norm R+\norm C$, integrate the other coordinate groups,
and apply the same argument to both lower-degree matrices.  Induction
therefore gives, even with the convenient larger linear constant $100$,
\[
             \mathbb E\norm Y
       \leq(200K)^q\mathcal L^{q/2}\sigma(A).
\]
Since $d=\max(d_1,d_2)$ and
\[
 \mathcal L\leq(q+4)\log(d+m),
\]
the exact advertised form follows with the explicit admissible constant
\[
                    C_q=200^q(q+4)^{q/2}.
\]
The constant is intentionally nonoptimal.  The proof preserves the
degree dependence, the exponent $q/2$, and the complete coefficient
flattening parameter of
\href{https://arxiv.org/html/2412.18468v2}{Theorem 2.4 of the matrix chaos paper}.
In particular $d=m=1$ is harmless because $\log(d+m)=\log2>0$.

\subsection{Strong chaos from the finite partial-transpose theorem}
\label{route:matrix-chaos-strong-khintchine}
\noindent\code{matrix-chaos-strong-khintchine}.

\begin{theorem}
Under the hypotheses of the two chaos norm cards, a stronger estimate is
\[
 \mathbb E\norm Y\leq(200K)^q
       \left\{\sigma(A)+\frac q2
          [\log(4d_1d_2m^q)]^{q/2}v(A)\right\}.
\]
In particular the requested bound with error exponent $(q+2)/2$ holds
with
\[
 C_q=200^q\left(1+
            \frac{q(q+4)^{q/2}}{2\log2}\right).
\]
\end{theorem}
\begin{proof}
First establish the strong linear inequality.  For rectangular
$B_j\in\R^{r\times c}$ write
\[
 \mathcal D(B_j)=\begin{pmatrix}0&B_j\\B_j^*&0\end{pmatrix},
 \qquad v_B=\norm{[\operatorname{vec}(B_j)^*]_j}.
\]
Then $\sigma(\mathcal D(B))=\sigma_B$.
As real vector spaces with Frobenius inner product,
$\mathcal D/\sqrt2$ is an isometry.  Applying it to each column of
the coefficient-vector matrix proves
$v(\mathcal D(B))=\sqrt2\,v_B$ exactly.
Theorem~\ref{thm:finite-hermite-partial-transpose} and the proved
$\kappa\leq v$ bound therefore yield
\[
 \mathbb E\norm{\sum_jg_jB_j}
 \leq2\sigma_B+6\sqrt2\,v_B\sqrt{\log(r+c)}.
\]
Lemma~\ref{lem:subgaussian-gaussian-domination}, followed by
$\sigma_B\leq\norm R+\norm C$, gives the convenient uniform version
\[
 \mathbb E\norm{\sum_jh_jB_j}
 \leq100K\left(\norm R+\norm C+\sqrt{\mathcal L}\,\norm V\right)
 \tag{SC}
\]
whenever $\mathcal L\geq\log(2(r+c))$.

Use the same fixed $\mathcal L=\log(4d_1d_2m^q)>1$ throughout the
recursion.  We prove the stronger claim with $r$ remaining groups and
the same global $\mathcal L$:
\[
 \mathbb E\norm{Y_r}
 \leq(200K)^r
       \{\sigma(Y_r)+(r/2)\mathcal L^{r/2}v(Y_r)\}.
\]
At $r=0$ this is the identity for a deterministic matrix.  At a positive
step apply (SC) conditionally on the first $r-1$ groups, which here means
an iterated integral on the finite product law.  Apply the strong
induction hypothesis to $R,C$, and the ordinary chaos estimate to $V$.
The finite bookkeeping lemma ensures that their variance and covariance
parameters are bounded by the indicated original parameters.  Hence,
writing $a=100$, the upper bound is at most
\begin{align*}
 aK(2aK)^{r-1}\{&2\sigma(A)
    +(r-1)\mathcal L^{(r-1)/2}v(A)
    +\mathcal L^{r/2}v(A)\}\\
 &\leq(2aK)^r\{\sigma(A)+(r/2)\mathcal L^{r/2}v(A)\}.
\end{align*}
The last step uses $\mathcal L\geq1$.  In the covariance branch the
frozen last index is a row coordinate, so its ultimate row subset is
nonempty; this is the reason the definition of $v$ excludes the empty
subset and still covers every term needed here.

Finally let $L_0=\log(d+m)\geq\log2$.  Use
$\mathcal L\leq(q+4)L_0$ and
$L_0^{q/2}\leq L_0^{(q+2)/2}/\log2$ to obtain the advertised conclusion
with the stated $C_q$.  Zero coefficients and one-dimensional matrices
are included without taking a logarithm of one.
\end{proof}

\paragraph{What changed relative to the original dependency graph.}
\href{https://arxiv.org/html/2412.18468v2}{Theorem 2.5 and Appendix A.3}
use a covariance-sensitive Gaussian norm theorem to supply their linear
step.  The later
\href{https://arxiv.org/html/2610.01982v1}{finite Hermite estimate}
supplies (SC) directly, with a square-root logarithm.  Thus the finite
proof above proves every displayed chaos bound, with explicit constants,
without assuming that deeper comparison theorem.  The covariance-sensitive
Gaussian norm card itself is also an immediate corollary: for $d\geq2$,
\[
 6\kappa\sqrt{\log d}
 \leq\frac{6}{(\log2)^{1/4}}
          \sqrt{\sigma v}\,(\log d)^{3/4},
\]
because $\kappa\leq\min(\sigma,v)\leq\sqrt{\sigma v}$.
This is a change of proof route, not a claim that the earlier sharper
free-model comparison has been proved by the same argument.


\section{Covariance bounds, free spectra, and block-Wishart variational formulas}

\subsection{Covariance-sensitive Gaussian norms without a free model}
\label{route:gaussian-series-covariance-sensitive-norm}
\noindent\code{gaussian-series-covariance-sensitive-norm}.
This card has a substantially shorter route than the original free-comparison
argument.  Use the finite Hermite theorem
\ref{thm:finite-hermite-partial-transpose} together with the coefficient-block
inequality that is now proved in the toolkit.

For Hermitian coefficients set
\[
 V=\sum_i A_i^2,\qquad
 C=\sum_i\operatorname{vec}(A_i)\operatorname{vec}(A_i)^*,\qquad
 T=(A_jA_i)_{i,j},\qquad
 \sigma=\norm V^{1/2},\quad v=\norm C^{1/2},\quad
 \kappa=\sqrt{\lambda_{\max}(T)}.
\]
All vectorizations here use the unnormalized Frobenius inner product.
The existing theorem
\code{NLAlib.coefficientBlock_variance_covariance_bounds}, verified with
\code{sorry=false}, gives
\[
 0\leq\kappa\leq\min\{\sigma,v\}\leq\sqrt{\sigma v}.
\]
Together with the finite Hermite norm estimate it follows that, for \(d\geq2\),
\begin{align*}
 \mathbb E\norm{\textstyle\sum_i g_iA_i}
 &\leq2\sigma+6\kappa\sqrt{\log d}\\
 &\leq2\sigma+
 \frac{6}{(\log2)^{1/4}}\sqrt{\sigma v}\,(\log d)^{3/4}.
\end{align*}
This proves the exact advertised leading constant \(2\), with an explicit
choice for the universal correction constant.  In fact the preceding
\(\kappa\sqrt{\log d}\) estimate is stronger.

For clarity, the finite mechanism behind the reused theorem is as follows.
On the finite space spanned by multiindices \(a\in\N^s\) with \(|a|\leq p\),
let \(C_i\) be truncated creation matrices and
\(L=\sum_i A_i\otimes C_i\).  On degrees at most \(p-1\),
\[
 L^*L=V\otimes I+B^*(T\otimes I)B,
 \qquad B^*B=\operatorname{diag}(|a|),
\]
where \(B\) stacks the annihilators \(C_i^*\).  On the top degree \(L=0\).
Hence \(\norm L^2\leq\sigma^2+(p-1)\kappa^2\).
The product Hermite recurrence identifies the squared norm of
\((L+L^*)^p\) applied to the vacuum with the Gaussian \(2p\)-th moment.
No infinite Fock space, asymptotic freeness, resolvent comparison, or
multilinear interpolation is needed for this card.
If the coefficient index set is empty, or \(\sigma=0\), every coefficient is
zero and the conclusion is immediate.  This also avoids any nonempty-index
side condition in the existing coefficient-block theorem.

\paragraph{Audit.}
The classical statement is
\href{https://arxiv.org/abs/2108.06312v2}{Bandeira--Boedihardjo--van Handel,
Corollary 2.2}.  The proof above is a different route through the finite
Hermite estimate.  It does \emph{not} prove comparison with the exact free
norm, which can be smaller than \(2\sigma\), or any reverse spectral inclusion.

\subsection{Two-sided free spectral comparison: finite geometric and probabilistic bridges}
\label{route:gaussian-free-spectrum-two-sided}
\noindent\code{gaussian-free-spectrum-two-sided}.
Assume \(d\geq2\). Write \(X=A_0+\sum_i g_iA_i\), with Hermitian coefficients, and let \(F=X_f\)
be the faithful free semicircular model.  Put
\[
 \delta=\sqrt{\sigma v},\qquad
 b=\sigma_*,\qquad \ell=\log d,
 \qquad R_Y(z)=(zI-Y)^{-1}.
\]
The following decomposition separates ordinary finite matrix analysis from
the two genuinely substantial free comparison inputs.  Unlike a direct
appeal to a spectral-set theorem, the inputs concern a single resolvent at
a single deterministic point.

\paragraph{Two exact analytic inputs still to prove.}
Let \(\tau_d=\operatorname{Tr}/d\) and use \(\tau_d\otimes\tau\) on the free
model.  For \(z=x+i\eta\), \(\eta>0\), \(q\geq1\), retain the resolvent
moment comparison
\begin{equation}
 \left|
 \bigl(\mathbb E\tau_d|R_X(z)|^{4q}\bigr)^{1/(4q)}
 -\norm{R_F(z)}_{L^{4q}(\tau_d\otimes\tau)}
 \right|
 \leq\frac83\frac{\delta^4(q+1)^3}{\eta^5}.
 \tag{RM}\label{eq:free-rm}
\end{equation}
The factor is obtained from the \(2p\)-resolvent statement with \(p=2q\).
The other input is the localized free resolvent estimate: for a universal
\(U\geq1\),
\begin{equation}
 \norm{R_F(x+i\eta)}
 \leq U\left(
 \norm{R_F(x+i\eta)}_{L^{4q}(\tau_d\otimes\tau)}
 +\frac b{\eta^2}\right),
 \quad x\in\operatorname{spec}F,\quad q\geq\log d.
 \tag{LU}\label{eq:free-lu}
\end{equation}
These are, respectively, the resolvent comparison in
\href{https://arxiv.org/abs/2108.06312v2}{BBvH, Theorem 6.1} and the
localized estimate in
\href{https://arxiv.org/abs/2406.11453v2}{BCSvH, Corollary 4.4}.
They are named prerequisites of this reroute, not consequences of the
covariance-sensitive norm theorem just proved.  In particular, an upper
norm bound alone cannot detect missing interior spectral points.

\paragraph{Parameter and Lipschitz identities.}
Finite Cauchy--Schwarz gives
\begin{equation}
 b\leq\min\{\sigma,v\}\leq\delta,
 \qquad \sigma\leq\sqrt d\,b.
 \label{eq:free-parameters}
\end{equation}
For the first inequality involving \(v\), test the covariance matrix against
the vectorization of a rank-one matrix of Frobenius norm one.  For the
inequality involving \(\sigma\), sum squared coordinates in an orthonormal
basis.  For the last inequality, fix a unit vector \(w\) and sum
\(\sum_i|e_j^*A_iw|^2\leq b^2\) over \(j\).
For real coefficient vectors \(g,h\),
\[
 \norm{X(g)-X(h)}\leq b\norm{g-h}_2,
 \qquad
 \bigl|\norm{R_{X(g)}(z)}-\norm{R_{X(h)}(z)}\bigr|
 \leq\frac b{\eta^2}\norm{g-h}_2.
\]
The second statement follows from the resolvent identity and
\(\norm{R_Y(x+i\eta)}\leq\eta^{-1}\).  Existing Gaussian Lipschitz
concentration therefore applies to the real-valued function
\(g\mapsto\norm{R_{X(g)}(z)}\).

\paragraph{Single-point bounds from the two inputs.}
Set \(q=\lceil\log d\rceil\), so \(q\geq1\) and
\(d^{1/(4q)}\leq e^{1/4}\).
Ordinary scalar Gaussian concentration and integration of its tail give
\[
 \bigl(\mathbb E\norm{R_X(z)}^{4q}\bigr)^{1/(4q)}
 \leq\mathbb E\norm{R_X(z)}+10\frac b{\eta^2}\sqrt q.
\]
For example, apply Minkowski to the mean and centered fluctuation, then
integrate the bound \(2e^{-r^2/(2L^2)}\), \(L=b/\eta^2\); the displayed
constant is deliberately generous.  The inequalities
\(\tau_d|R|^{4q}\leq\norm R^{4q}\leq d\tau_d|R|^{4q}\),
\eqref{eq:free-rm}, and concentration now give, with probability at least
\(1-e^{-r^2}\),
\begin{equation}
 \norm{R_X(z)}
 \leq e^{1/4}\norm{R_F(z)}
  +e^{1/4}\frac83\frac{\delta^4(q+1)^3}{\eta^5}
  +\sqrt2\frac b{\eta^2}r.
 \label{eq:free-point-upper}
\end{equation}
For \(x\in\operatorname{spec}F\), \eqref{eq:free-lu} and the reverse
direction of \eqref{eq:free-rm} give, with the same probability,
\begin{equation}
 \norm{R_X(x+i\eta)}
 \geq\frac1{U\eta}
  -\frac83\frac{\delta^4(q+1)^3}{\eta^5}
  -\frac b{\eta^2}(1+10\sqrt q+\sqrt2r).
 \label{eq:free-point-lower}
\end{equation}
Here \(\norm{R_F(x+i\eta)}=1/\eta\) because \(x\) belongs to its spectrum.
These estimates make the separate roles of moment comparison and
localized ultracontractivity explicit.

\paragraph{Finite grids, with no dependence on the size of the mean.}
The deterministic free norm bound is
\(\norm{F-A_0\otimes1}\leq2\sigma\).
Consequently \(\operatorname{spec}F\) lies in the union of \(d\) intervals
of radius \(2\sigma\) about the eigenvalues of \(A_0\).  By
\eqref{eq:free-parameters}, this set has a \(b\)-net whose points belong to
\(\operatorname{spec}F\) and whose cardinality is at most \(6d^{3/2}\).
Take a maximal \(b\)-separated subset within each interval and count the
disjoint gaps; no interval between widely separated eigenvalues of \(A_0\)
needs to be covered.

For the other direction, the classical Gaussian-series bound and Gaussian
concentration imply, except on an event of probability \(e^{-r^2}\),
\[
 \norm{X-A_0}\leq
 \sigma\sqrt{2\log(2d)}+\sqrt2\,b r=:M_r.
\]
Thus the random spectrum lies in the union of the same \(d\) intervals,
now of radius \(M_r\).  A fixed grid of mesh at most \(b\) for these intervals
has at most \(8d^2(1+r)\) points.  Here one can use
\(\log(2d)\leq d\) for \(d\geq2\).
This coarse polynomial bound suffices.  Its logarithm is bounded by a
universal multiple of \(\log d+r^2\).

Use \eqref{eq:free-point-upper} on this second grid, and
\eqref{eq:free-point-lower} on the first.  More explicitly, for fixed
\(t\geq0\), set
\[
 r=\sqrt{t^2+\log3},\qquad N_0=8d^2(1+r),\qquad
 a=\sqrt{t^2+\log(3N_0)}.
\]
The norm-envelope failure probability is \(e^{-t^2}/3\).
Both grids have at most \(N_0\) points, and using pointwise tail parameter
\(a\) makes each grid failure probability at most \(e^{-t^2}/3\).
These grids are deterministic: their radii depend on \(t\), not on the
sample.  To bound the resulting error, note that \(r\leq t+2\) and
\(t+3\leq4e^{t^2/2}\), so
\[
 \log(3N_0)\leq\log96+2\ell+t^2/2
                 \leq9\ell+t^2/2,
 \qquad a\leq2t+3\sqrt\ell.
\]
The scalar resolvent identity
\[
 \bigl|\norm{R_Y(x+i\eta)}-\norm{R_Y(x'+i\eta)}\bigr|
 \leq\frac{|x-x'|}{\eta^2}
\]
extends the lower estimate with an additional \(b/\eta^2\) error and the
upper estimate with \((1+e^{1/4})b/\eta^2\).  In the upper estimate both
resolvents change with the grid point.
Using this bound for \(a\), \(q\leq(5/2)\ell\),
\(q+1\leq4\ell\), and \(e^{1/4}<4/3\), all errors can be bounded by
\[
 250\frac{\delta^4\ell^3}{\eta^5}
 +40\frac{b(\sqrt\ell+t)}{\eta^2}.
\]
The harmless constants for \(\ell\in[\log2,1]\) must be retained:
\(q+1\leq(1+2/\log2)\ell\), and
\(b\sqrt\ell\leq(\log2)^{-1/4}\delta\ell^{3/4}\).

Choose
\[
 \eta=D\bigl(\delta\ell^{3/4}+bt\bigr)
\]
with \(D=1000U\): since
\(b\sqrt\ell\leq2\delta\ell^{3/4}\), the error multiplied by \(\eta\)
is at most \(250/D^4+80/D\leq1/(2U)\).
The errors are therefore at most
\(\min\{1/2,1/(2U)\}/\eta\).
At a random spectral point \(x\in\operatorname{spec}X\), the left side of
\eqref{eq:free-point-upper} equals \(1/\eta\).  It follows that
\(\norm{R_F(x+i\eta)}\geq c/\eta\), and hence
\(\operatorname{dist}(x,\operatorname{spec}F)\leq3\eta\),
using \(2e^{1/4}<3\).
At a free spectral point, \eqref{eq:free-point-lower} similarly gives
\(\operatorname{dist}(x,\operatorname{spec}X)\leq2U\eta\).
We have proved, conditional exactly on (RM) and (LU),
\[
 \mathbb P\left\{d_H(\operatorname{spec}X,\operatorname{spec}F)
 >C\delta(\log d)^{3/4}+Cbt\right\}\leq e^{-t^2},
 \qquad C=3000U^2.
\]
No arbitrary compact-set entropy theorem is required: only finite unions of
intervals and a finite union bound.  If \(b=0\), all coefficients vanish
and both spectral sets equal \(\operatorname{spec}A_0\); handle this before
constructing any grid or choosing \(\eta\).

For nonempty compact subsets of \(\R\), the maps minimum, maximum, and
maximum absolute value are all 1-Lipschitz for Hausdorff distance.  Integrating
the displayed tail yields each expectation assertion, since
\[
 \int_0^\infty e^{-t^2}\,dt=\frac{\sqrt\pi}{2},
 \qquad b\leq\delta,
 \qquad (\log d)^{3/4}\geq(\log2)^{3/4}.
\]
For example, \(9000U^2\) is a single constant valid for both the stated
tail and all three expectation inequalities: integration multiplies the
tail constant by at most
\(1+\sqrt\pi/[2(\log2)^{3/4}]<3\).
The theorem controls sets, not eigenvalue multiplicities or ordered matching.

\paragraph{A concrete free model and a finite algebra interface.}
There is no need to develop an abstract theory of every noncommutative
probability space first.  On the Hilbert space with orthonormal basis all
finite words over \(\{1,\ldots,s\}\), let \(L_i e_w=e_{iw}\) and
\(s_i=L_i+L_i^*\).  Left creation satisfies \(L_i^*L_j=\delta_{ij}I\), so
\[
 \norm{\sum_i A_i\otimes L_i}^2=\norm{\sum_iA_i^2},
 \qquad \norm{\sum_iA_i\otimes s_i}\leq2\sigma.
\]
The vacuum moment of a word in the \(s_i\)'s is its finite noncrossing-pairing
sum: expansion into creation and annihilation operators proves this by
induction on word length.  These moments are cyclic, so the vacuum state is
tracial.  It is faithful on the generated algebra: right fields commute with
left fields, their vacuum polynomials span every finite word, and an element
annihilating the vacuum consequently annihilates a dense subspace.  The same
argument applies to matrices of such elements.

Every vacuum moment of order at most \(2p\) can be computed in the finite
word space of lengths at most \(p\).  This gives a concrete target for the
combinatorial parts of (RM).  It does not justify replacing the operator
spectrum by the spectrum of a single finite truncation: truncation creates
boundary eigenvalues.  The uniform resolvent moment comparison (RM) and the
localized bound (LU), including their constants independent of \(A_0\),
remain the analytic work for this card.

\subsection{Quadratic spectra from one shifted block matrix}
\label{route:gaussian-quadratic-free-spectrum}
\noindent\code{gaussian-quadratic-free-spectrum}.
For \(d\geq2\), this card reduces completely to the preceding affine spectral comparison.
All additional work is a deterministic block identity and scalar estimates;
there is no new nonlinear Gaussian comparison theorem.

\begin{lemma}[Shifted Gram linearization, with both spectral directions]
Let \(X,Y\) be square matrices or bounded operators in their respective
unital \(C^*\)-algebras, let \(B=B^*\), and let \(\varepsilon>0\).  Form
\[
 B_\varepsilon=B+(\norm B+4\varepsilon^2)I,\qquad
 H_\varepsilon(X)=
 \begin{bmatrix}0&X&B_\varepsilon^{1/2}\\
 X^*&0&0\\B_\varepsilon^{1/2}&0&0\end{bmatrix}.
\]
If the spectra of \(H_\varepsilon(X)\) and \(H_\varepsilon(Y)\) have
Hausdorff distance at most \(\varepsilon\), then
\[
 d_H\bigl(\operatorname{spec}(XX^*+B),
          \operatorname{spec}(YY^*+B)\bigr)
 \leq4\bigl(\norm Y+\sqrt{\norm B}\bigr)\varepsilon+5\varepsilon^2.
\]
\end{lemma}
\begin{proof}
Write \(Z=[X\ B_\varepsilon^{1/2}]\).  The displayed block matrix is the
Hermitian dilation of \(Z\), and \(ZZ^*=XX^*+B_\varepsilon\succeq
4\varepsilon^2I\).  Its nonzero spectrum consists of the positive and
negative square roots of \(\operatorname{spec}(ZZ^*)\).  This follows from
the block equations \(Zw=\lambda u\), \(Z^*u=\lambda w\), or from the block
inverse when \(\lambda\neq0\); it works for operators as well as matrices.

If \(a\) is a positive spectral square root on one side, then
\(a\geq2\varepsilon\).  A spectral point on the other side at distance at
most \(\varepsilon\) is strictly positive, so it cannot be an artificial
zero or a negative square root.  Therefore the positive square-root spectral
sets themselves have Hausdorff distance at most \(\varepsilon\).
Let
\(M=\norm{YY^*+B_\varepsilon}^{1/2}\).
For a point \(a\) from the \(X\)-side choose \(b\) from the \(Y\)-side with
\(|a-b|\leq\varepsilon\); then \(b\leq M\) and
\(|a^2-b^2|\leq2M\varepsilon+\varepsilon^2\).
For the reverse direction start with \(b\leq M\) and use exactly the same
inequality.  In particular, the reverse estimate does not introduce
\(\norm X\).
Translate both squared spectra by \(-\norm B-4\varepsilon^2\), and use
\[
 M\leq\sqrt{\norm Y^2+2\norm B+4\varepsilon^2}
 \leq\norm Y+\sqrt{2\norm B}+2\varepsilon.
\]
The claimed, slightly coarser, constants follow.
\end{proof}

For \(X=A_0+\sum_i g_iA_i\), the random coefficients in the linearization
are
\[
 D_i=\begin{bmatrix}0&A_i&0\\A_i^*&0&0\\0&0&0\end{bmatrix}.
\]
Direct block multiplication gives
\[
 \sigma(D)=\sigma(A),\qquad v(D)\leq\sqrt2\,v(A),\qquad
 \sigma_*(D)=\sigma_*(A),\qquad
 \delta(D)\leq2^{1/4}\delta(A).
\]
For the covariance inequality split a vector against the upper and lower
off-diagonal vectorizations and use \(|a+b|^2\leq2|a|^2+2|b|^2\).
For the weak variance upper bound use
\[
 \left(\sum_i|u_1^*A_iw_2+u_2^*A_i^*w_1|^2\right)^{1/2}
 \leq\sigma_*(A)
  (\norm{u_1}\norm{w_2}+\norm{u_2}\norm{w_1})
 \leq\sigma_*(A)\norm u\norm w.
\]
The reverse inequality follows by restricting to \(u_2=w_1=0\).

If the affine comparison constant in the preceding subsection is \(C_a\),
define the numerical factor
\[
 \gamma=2^{1/4}
 \left(1+\frac{\log3}{\log2}\right)^{3/4},
 \qquad D=C_a(\gamma+1).
\]
For \(t\geq\delta(\log d)^{3/4}\), \(\sigma_*>0\), use the affine theorem
in dimension \(3d\) with tail parameter \(t/\sigma_*\) and choose
\(\varepsilon=Dt\).  Since \(\log(3d)\leq
(1+\log3/\log2)\log d\), the resulting exceptional probability is at most
\(\exp(-t^2/\sigma_*^2)\).  The deterministic lemma then gives the exact
form in the card, with any universal
\(C\geq\max\{4D,5D^2\}\).
If \(\sigma_*=0\), all \(A_i=0\), and both original spectral sets agree
without a tail argument.

\paragraph{Reuse and audit.}
The verified declarations \code{NLAlib.dilation},
\code{NLAlib.dilation_sq}, and \code{NLAlib.dilation_identities} provide the
finite matrix algebra.  The operator version needs the same block inverse,
not a new finite-dimensional SVD.  The positive shift is essential for the
reverse inclusion.  The dimension in the affine theorem is \(3d\), not
\(d\); the factor \(\gamma\) above records the resulting logarithm.  This
matches the mechanism of
\href{https://arxiv.org/abs/2406.11453v2}{BCSvH, Lemma 5.3}, with the
dimension conversion made explicit.

\subsection{Block-Wishart edges: simplify the finite part of global inversion}
\label{route:block-wishart-variational-edges}
\noindent\code{block-wishart-variational-edges}.
Let \(k\geq1\), \(\alpha>0\), and let \(\nu\) be a compactly supported
probability law on Hermitian \(k\times k\) matrices.  Throughout,
\(\operatorname{Tr}\) is unnormalized.  Define
\[
 K(S)=\int T(I+ST)^{-1}\,d\nu(T)-\alpha^{-1}S^{-1},
 \qquad
 \mathcal D_0=\{S\succ0:S^{-1}+T\succ0\text{ on }\operatorname{supp}\nu\}.
\]
Write \(\mathcal P_-\) for the points of \(\mathcal D_0\) at which \(DK\)
is positive definite as an operator on the real Hilbert space of Hermitian
matrices with Frobenius inner product.

\paragraph{The now-proved local interface.}
Set
\[
 C_S(T)=T(I+ST)^{-1}=(I+TS)^{-1}T,
 \quad W_S(T)=S^{1/2}C_S(T)S^{1/2},
 \quad \mathcal L_S(H)=\alpha\int W_S(T)H W_S(T)\,d\nu(T).
\]
These matrices \(C_S(T),W_S(T)\) are Hermitian even for indefinite \(T\).
For \(H=S^{-1/2}\Delta S^{-1/2}\),
\[
 DK(S)[\Delta]=\alpha^{-1}S^{-1/2}
              (H-\mathcal L_S(H))S^{-1/2}.
\]
The following statements have already been verified with \code{sorry=false}:
\code{NLAlib.hasFDerivAt_matrixKTransform_of_posDef},
\code{NLAlib.matrixKTransformDeriv_eq_congruence},
\code{NLAlib.matrixKTransform_stability_iff}, and
\code{NLAlib.matrixKTransform_supersolution_iff_derivative_witness}.
In particular,
\begin{equation}
 S\in\mathcal P_-
 \quad\Longleftrightarrow\quad
 S\in\mathcal D_0\ \text{and there is }B\succ0
 \text{ with }\mathcal L_S(B)\prec B.
 \label{eq:wishart-witness}
\end{equation}
The map \(\mathcal L_S\) preserves positive semidefinite matrices and is
self-adjoint in Frobenius geometry.  It need not be positive semidefinite
in that Hilbert-space geometry.  For example, conjugation by
\(\operatorname{diag}(1,-1)\) negates the Hermitian off-diagonal matrix.
Thus the erroneous stronger assertion in the source's discussion following
its equation (103) must not be reused.  The verified stability equivalence
above already supplies the corrected argument.

\begin{lemma}[A single weighted contraction proves global injectivity]
The map \(K\) is injective on \(\mathcal P_-\).
\end{lemma}
\begin{proof}
Fix \(S,R\in\mathcal P_-\).  The inverse identity gives
\[
 \alpha S^{1/2}(K(S)-K(R))R^{1/2}
   =(I-\mathcal G)H,
 \quad H=S^{-1/2}(S-R)R^{-1/2},
 \quad\mathcal G(H)=\alpha\int W_S(T)H W_R(T)\,d\nu(T).
\]
Choose witnesses \(B,D\succ0\) in \eqref{eq:wishart-witness} and numbers
\(0\leq a,b<1\) satisfying
\[
 \mathcal L_S(B)\preceq aB,\qquad
 \mathcal L_R(D)\preceq bD.
\]
These numbers can be the largest eigenvalues of the corresponding
positive congruences.  Conjugate \(\mathcal G\) by
\(H=B^{1/2}ZD^{1/2}\).  The resulting map is
\[
 \widetilde{\mathcal G}(Z)=\alpha\int P(T)ZQ(T)\,d\nu(T),
 \quad P=B^{-1/2}W_SB^{1/2},\quad Q=D^{1/2}W_RD^{-1/2}.
\]
For unit vectors \(u,v\), scalar integral Cauchy--Schwarz yields
\begin{align*}
 |u^*\widetilde{\mathcal G}(Z)v|
 &\leq\norm Z
 \left(\alpha\int\norm{P^*u}^2d\nu\right)^{1/2}
 \left(\alpha\int\norm{Qv}^2d\nu\right)^{1/2}\\
 &\leq\sqrt{ab}\,\norm Z.
\end{align*}
Consequently \(\norm{\widetilde{\mathcal G}}\leq\sqrt{ab}<1\) in the
ordinary matrix operator norm.  Thus \(I-\mathcal G\) is invertible by a
Neumann series in this equivalent weighted norm.  If \(K(S)=K(R)\), the
displayed difference identity forces \(H=0\), hence \(S=R\).
\end{proof}
This replaces products over \(2p\) copies of the measure and estimates of
iterated Frobenius norms by one Cauchy--Schwarz inequality.  It does not need
\(\mathcal G\) to preserve Hermitian matrices.

\begin{lemma}[An explicit path to a common convex region]
The stable domain \(\mathcal P_-\) is open and path connected.
\end{lemma}
\begin{proof}
Openness follows from the uniform positive gap on the compact support,
continuity of inversion, and strict positivity of the finite-dimensional
derivative.  For connectedness set \(Y=S^{-1}\) and write the witness
condition in the equivalent form
\[
 \mathcal F_Y(B_0):=
 \alpha\int T(Y+T)^{-1}B_0(Y+T)^{-1}T\,d\nu(T)\prec B_0,
 \quad B_0\succ0.
\]
Here $B_0=Y^{1/2}BY^{1/2}$ converts a witness for $\mathcal L_S$ into
a witness for $\mathcal F_Y$; congruence by $Y^{1/2}$ verifies the equivalence.
If this holds, it continues to hold on the whole ray \(Y_t=Y+tB_0\),
\(t\geq0\).  Indeed, put
\(Q=B_0^{-1/2}(Y+T)B_0^{-1/2}\succ0\).  Functional calculus of this single
matrix gives \((Q+tI)^{-2}\preceq Q^{-2}\), hence
\[
 (Y_t+T)^{-1}B_0(Y_t+T)^{-1}
 \preceq(Y+T)^{-1}B_0(Y+T)^{-1}.
\]
Congruence by \(T\) and integration preserve this inequality.  This is a
commuting functional-calculus argument; it does not assert that the
function \(x\mapsto x^{-2}\) is operator monotone.

Let \(M=\max_{T\in\operatorname{supp}\nu}\norm T\) and
\(c=\sqrt\alpha M+1\).  The region
\[
 \mathcal G_0=\{Y\succ(M+c)I\}
\]
is convex and lies in the stable domain in inverse coordinates, since
\(\mathcal F_Y(I)\preceq\alpha M^2c^{-2}I\prec I\).
Every ray \(Y+tB_0\) eventually enters this region; for example take
\(t>(M+c-\lambda_{\min}(Y))_+/\lambda_{\min}(B_0)\).
Two rays and a straight segment inside \(\mathcal G_0\) join any two points.
Invert this path to obtain a path in \(\mathcal P_-\).
\end{proof}

\paragraph{The precise remaining global bridge.}
The two lemmas above and the inverse function theorem make \(K\) a
diffeomorphism from \(\mathcal P_-\) onto its open image \(\mathcal U\).
They do \emph{not} identify this image with the correct spectral branch.
For the characterized probability law in the card, the remaining bridge is
\begin{equation}
 \mathcal U=\{U=U^*:W-U\otimes1\succ0\},\qquad
 \operatorname{spec}W=\operatorname{supp}\mu_{\alpha,\nu},
 \label{eq:wishart-global-bridge}
\end{equation}
for a faithful bounded realization \(W\), with
\(K^{-1}(uI)=S(u)\) on the scalar line.  The relevant normalization is
\[
 S(u)=\alpha^{-1}(\operatorname{id}\otimes\tau)
                (W-uI)^{-1},\qquad
 \frac\alpha k\operatorname{Tr}S(u)
       =\int\frac{d\mu_{\alpha,\nu}(\lambda)}{\lambda-u}.
\]
This is the substantive branch-identification part of
\href{https://arxiv.org/abs/2606.27774v1}{Montanari--Saeed, Theorem 3}.
Its proof requires more than local derivative positivity.  In particular,
neither connectedness nor global injectivity alone implies surjectivity
onto the spectral positivity domain.

\paragraph{A real-variable replacement for the final continuation argument.}
There is a useful further reduction of the global bridge.  Suppose a
faithful bounded realization \(W\) has been constructed, and its outside
resolvent transform \(\widehat S(U)\) has been proved to lie in
\(\mathcal P_-\) and satisfy \(K(\widehat S(U))=U\) for
\(U\in\mathcal O:=\{U:W-U\otimes1\succ0\}\).
These are still prerequisites; the scalar characterization in the card
does not by itself establish them for nonscalar \(U\).
The remaining surjectivity then follows without contour integration or a
general Stieltjes inversion theorem.

Here is the elementary scalar obstruction needed.  If \(\rho\) is a
compactly supported positive measure on \([0,\infty)\), with
\(0\in\operatorname{supp}\rho\), then
\[
 h(x)=\int(\lambda-x)^{-1}\,d\rho(\lambda),\qquad x<0,
\]
cannot have a real-analytic continuation to a neighborhood of zero.
Indeed, such a continuation would have Taylor coefficients bounded by
\(C R^{-n}\) for some \(C,R>0\).  Differentiation at negative \(x\),
followed by monotone convergence as \(x\uparrow0\), would give
\[
 \int\lambda^{-n-1}\,d\rho(\lambda)
   =\frac{h^{(n)}(0)}{n!}\leq C R^{-n}.
\]
For any \(0<r<R\), support at zero gives \(\rho([0,r])>0\); the left
side is at least \(\rho([0,r])r^{-n-1}\), contradicting the bound as
\(n\to\infty\).  An atom at zero already contradicts the \(n=0\) bound.

Apply this obstruction as follows.  The open set \(\mathcal O\) is
contained in the connected open image \(K(\mathcal P_-)\).  If the
containment were proper, a path in the latter image would meet
\(U_0\in\partial\mathcal O\).  Here
\(W-U_0\otimes1\succeq0\), and zero belongs to its spectrum and, by
faithfulness, to its scalar spectral support.  The local real-analytic
inverse of \(K\) at \(U_0\) agrees with \(\widehat S\) just inside
\(\mathcal O\), by global injectivity.  Its scalar trace on
\(U_0+xI\), \(x<0\), would therefore analytically continue the above
\(h\) across zero, a contradiction.  The analyticity needed here is only
that of the finite matrix map \(K\) and its local inverse; uniform
Neumann expansions on the compact \(T\)-support give analyticity of
\(K\).  This closes the final global-coverage step once the stated
outside-resolvent prerequisites are available.  It does not assume that
an arbitrary complex continuation across a spectral cut is impossible:
the real continuation from below a positive spectral endpoint is the
property used.

Once \eqref{eq:wishart-global-bridge} is supplied, the edge formula is a
short order argument.  If \(U\in\mathcal U\), then
\(W\succ U\otimes1\succeq\lambda_{\min}(U)I\), so
\(\lambda_{\min}(U)\leq\inf\operatorname{spec}W\).
Conversely, every \(u<\inf\operatorname{spec}W\) has \(uI\in\mathcal U\).
Taking suprema and changing coordinates through \(U=K(S)\) gives
\[
 \inf\operatorname{supp}\mu_{\alpha,\nu}
   =\sup_{S\in\mathcal P_-}\lambda_{\min}(K(S)).
\]
Strict positivity of the left endpoint is equivalent to some \(S\in
\mathcal P_-\) having \(K(S)\succ0\); substituting the already-proved
derivative-witness equivalence gives the exact \(S,B\succ0\) certificate
in the card.  The upper edge follows from reflection \(T\mapsto-T\).
No convergence of the extreme eigenvalues of finite random samples is
asserted by this argument.

\subsection{Block-Wishart log potential through a finite inverse Hessian}
\label{route:block-wishart-log-potential}
\noindent\code{block-wishart-log-potential}.
Retain the notation above and the explicit global bridge
\eqref{eq:wishart-global-bridge}.  Only the following consequences of that
bridge are needed here: \(\mathcal U=K(\mathcal P_-)\) is convex;
\(uI\in\mathcal U\) exactly for \(u<\zeta_-\); and its scalar inverse branch
has the displayed Stieltjes trace.  The minimization proof itself can be
carried out entirely in the \(k^2\)-dimensional real space of Hermitian
matrices.

For \(S\in\mathcal P_-\), define
\[
 g(S;u)=-\alpha u\operatorname{Tr}S
       +\alpha\int\log\det(I+TS)\,d\nu(T)
       -\log\det S-k(\log\alpha+1).
\]
There is a convenient all-Hermitian spelling of the logarithm:
\begin{equation}
 \log\det(I+TS)=\log\det S+\log\det(S^{-1}+T).
 \label{eq:wishart-log-hermitian}
\end{equation}
Both determinants on the right are strictly positive.  Equivalently,
\(I+TS\) is similar to \(I+S^{1/2}TS^{1/2}\succ0\).  Thus there is no
complex logarithm and no choice of branch.  Compactness of the support
gives a uniform positive spectral gap locally in \(S\), justifying
differentiation under this finite matrix-valued integral.

Jacobi's determinant identity and the inverse derivative give
\begin{equation}
 D_Sg(S;u)[H]=\alpha\operatorname{Tr}((K(S)-uI)H),
 \qquad D_S^2g(S;u)[H,H]=\alpha\langle H,DK(S)[H]\rangle_F.
 \label{eq:wishart-gradient}
\end{equation}
For a direct proof of the first identity, differentiate
\(\log\det(I+TS)\) to get
\(\operatorname{Tr}((I+TS)^{-1}TH)=\operatorname{Tr}(C_S(T)H)\), and
differentiate \(\log\det S\) to get \(\operatorname{Tr}(S^{-1}H)\).
The traces are real because their factors are Hermitian.
The finite determinant derivative can itself be proved by expanding
\(\det(I+tA)=1+t\operatorname{Tr}A+O(t^2)\); this uses no complex analysis.
The existing real-matrix theorem \code{NLAlib.hasDerivAt_det} is available,
but its complex-Hermitian transport must be supplied explicitly if used.

\begin{theorem}[Global minimum assuming the branch bridge]
For every \(u<\zeta_-\), the unique minimum of \(g(\cdot;u)\) on
\(\mathcal P_-\) occurs at \(S_u=K^{-1}(uI)\).
\end{theorem}
\begin{proof}
Fix any \(S_1\in\mathcal P_-\), put \(\Delta=K(S_1)-uI\), and use the
straight segment \(U_t=uI+t\Delta\) in the convex image \(\mathcal U\).
Let \(S_t=K^{-1}(U_t)\).  Ordinary finite-dimensional inverse differentiation
gives
\[
 \dot S_t=DK(S_t)^{-1}[\Delta].
\]
By \eqref{eq:wishart-gradient},
\[
 \frac d{dt}g(S_t;u)
  =\alpha t\langle\Delta,DK(S_t)^{-1}[\Delta]\rangle_F\geq0.
\]
The inverse of a positive definite self-adjoint operator is positive
definite, so this derivative is strictly positive for \(t>0\) when
\(\Delta\neq0\).  Integrating proves the minimum and uniqueness.
This argument does not require \(\mathcal P_-\) to be convex.  Positivity
of the Hessian on a nonconvex domain, by itself, would not prove a global
minimum.
\end{proof}

It remains to identify the minimum value, including its constant.
Differentiating along the critical scalar branch cancels the \(S_u\)
derivative and gives
\[
 \frac d{du}\frac1k g(S_u;u)
  =-\frac\alpha k\operatorname{Tr}S_u
  =-\int\frac{d\mu_{\alpha,\nu}(\lambda)}{\lambda-u}
  =\frac d{du}\int\log(\lambda-u)\,d\mu_{\alpha,\nu}(\lambda).
\]
The difference is constant on \(( -\infty,\zeta_-)\).
To determine it, set \(u=-r\) and write
\(S_{-r}=(\alpha r)^{-1}R_r\).  The branch at infinity satisfies
\(R_r=I+O(r^{-1})\).  This expansion can be obtained directly from the
finite equation
\[
 R_r^{-1}=I+\frac1r\int
 T\bigl(I+(\alpha r)^{-1}R_rT\bigr)^{-1}\,d\nu(T)
\]
by the finite-dimensional implicit function theorem at \(r^{-1}=0,R=I\);
the global bridge identifies this local solution with the Stieltjes branch.
Consequently,
\[
 -\alpha(-r)\operatorname{Tr}S_{-r}=k+o(1),\quad
 \int\log\det(I+TS_{-r})\,d\nu=o(1),\quad
 -\log\det S_{-r}=k\log(\alpha r)+o(1).
\]
It follows that \(k^{-1}g(S_{-r};-r)-\log r\to0\).
Compact support also gives
\(\int\log(\lambda+r)d\mu-\log r\to0\), so the constant is zero.  Thus
\[
 \int\log(\lambda-u)\,d\mu_{\alpha,\nu}(\lambda)
   =\frac1k\min_{S\in\mathcal P_-}g(S;u).
\]
The factor \(\alpha/k\) in the Stieltjes trace and the term
\(-k(\log\alpha+1)\) are both necessary; dropping either changes this
normalization.  Reflection gives the analogous formula to the right of
the support.

\paragraph{Status of the reroutes.}
The covariance norm card is fully reduced to finite Hermite algebra and
already-proved coefficient-block bounds.  The quadratic card is a complete
reduction to the affine spectral comparison.  For the affine card, (RM)
and (LU) remain explicit analytic prerequisites; the elementary
concentration, finite-grid, and Hausdorff steps above are proved.
For the two block-Wishart cards, global branch identification
\eqref{eq:wishart-global-bridge} remains a prerequisite.  The weighted
contraction, connectedness, edge optimization, minimum, and normalization
arguments above are independent finite-dimensional derivations around
that one core.  Weak convergence of sample laws alone does not justify
either an edge limit or integration of an unbounded logarithm; those
finite-sample assertions are not part of these cards.


\section{Gaussian replacement, exchangeable covariance, and resolvent comparison}

All matrices in this section are finite-dimensional.  The Hermitian matrices
are regarded as a real vector space, including in the complex case.  Thus
``matching covariance'' always means equality of every real-coordinate
second moment.  Equality of $\mathbb E H^2$ alone is insufficient.
Independent copies and Gaussian proxies can be constructed on a product
of their finite-dimensional laws.  Every asserted probability or expectation
then transfers back by equality in distribution.  No standard-Borel
assumption on the original probability space is needed.

\subsection{Gaussian replacement for the trace MGF}
\label{route:gaussian-trace-mgf-comparison}
\code{gaussian-trace-mgf-comparison}.
The reroute uses one quadratic form, scalar Taylor's formula, and symmetry.
It retains the deterministic BMV/Stahl theorem as an explicit deep input;
the argument below does not prove that theorem.

\paragraph{The precise retained input.}
For Hermitian $A,H$, write $F_H(t)=\operatorname{Tr}e^{A+tH}$.
BMV/Stahl supplies a positive finite measure $\nu_{A,H}$ supported on
$[\lambda_{\min}(H),\lambda_{\max}(H)]$ such that
\[
 F_H(t)=\int e^{tx}\,d\nu_{A,H}(x)\qquad(t\in\R).
\]
Consequently, for $H\preceq RI$ and $0\le t\le1$,
\begin{equation}\label{eq:comparison-bmv-two}
 0\le F_H''(t)\le e^{tR}F_H''(0),\qquad
 \frac{F_H(a)+F_H(-a)}2\ge F_H(0)+\frac{a^2}2F_H''(0).
\end{equation}
The first inequality follows by integrating $x^2e^{tx}$; the second follows
from $\cosh(ax)\ge1+a^2x^2/2$.  These two scalar consequences are the entire
BMV interface required here.  The reference is
\href{https://arxiv.org/abs/1107.4875v3}{Stahl's BMV theorem}; the comparison
target is \href{https://arxiv.org/html/2603.04365v5#S4.SS6}{Tropp,
Propositions 4.6--4.8}.  In particular, no random-offset conditional Lieb
inequality is involved in this replacement argument.

\paragraph{A finite covariance calculation.}
The first two directional derivatives are
\[
 F_H'(0)=\operatorname{Tr}(He^A),\qquad
 Q_A(H):=F_H''(0)
 =\int_0^1\operatorname{Tr}\bigl(H e^{(1-s)A}H e^{sA}\bigr)\,ds.
\]
The second identity follows by differentiating the exponential once and
using cyclicity of the trace.  In an eigenbasis $A=\diag(a_1,\ldots,a_d)$,
\[
 Q_A(H)=\sum_{i,j}|H_{ij}|^2
       \int_0^1e^{(1-s)a_j+sa_i}\,ds.
\]
This is a nonnegative quadratic form on the real coordinates of $H$,
bounded by $d e^{\norm{A}}\norm{H}^2$.  Thus finite second moments suffice
for its integrability, and covariance matching gives
$\mathbb E Q_A(W)=\mathbb E Q_A(G)$.

\paragraph{One replacement, with only an upper bound.}
Assume $\mathbb EW=0$, $W\preceq RI$, $R\ge0$, and
$\mathbb E\norm{W}^2<\infty$.  Let $G$ be centered Gaussian with the full
covariance of $W$.  Set
\[
 \gamma_R^2=2\int_0^1(1-s)e^{Rs}\,ds.
\]
Apply Taylor's formula pointwise to $F_W$ on $[0,1]$, and then the first
inequality in \eqref{eq:comparison-bmv-two}.  Integration gives
\[
 \mathbb E\operatorname{Tr}e^{A+W}
 \le\operatorname{Tr}e^A+\frac{\gamma_R^2}{2}\mathbb E Q_A(W).
\]
Since $G$ and $-G$ have the same law, the second inequality in
\eqref{eq:comparison-bmv-two} gives
\[
 \mathbb E\operatorname{Tr}e^{A+\gamma_RG}
 =\mathbb E\frac{F_G(\gamma_R)+F_G(-\gamma_R)}2
 \ge\operatorname{Tr}e^A+\frac{\gamma_R^2}{2}\mathbb E Q_A(G).
\]
This proves the one-step comparison.  It avoids differentiating an
expectation on an interval of negative parameters: such expectations
need not exist for an upper-bounded, heavy-negative-tailed $W$.
The left side is bounded by $d e^{\lambda_{\max}(A)+R}$; Gaussian
exponential moments make the right side finite.

\paragraph{Independent sums and all endpoints.}
Let $W_i$ be independent centered Hermitian matrices with finite second
moments and $W_i\preceq RI$.  On a product space choose independent
centered $G_i$ with the covariance of $W_i$, also independent of all $W_i$.
Apply the preceding deterministic-offset inequality with
$W\leftarrow\theta W_i$ and offset equal to the sum of all other terms,
then integrate over that offset.  Finitely many replacements give, for
$\theta\ge0$,
\begin{equation}\label{eq:comparison-trace-mgf}
 \mathbb E\operatorname{Tr}e^{\theta(\sum_iW_i+\Delta)}
 \le\mathbb E\operatorname{Tr}e^{g_R(\theta)Z+\theta\Delta},
 \quad
 g_R(\theta)=\theta\sqrt{2\int_0^1(1-s)e^{\theta Rs}\,ds},
\end{equation}
where $Z$ has the law of $\sum_iG_i$.  The mixed intermediate expectations
are finite: the unreplaced summands have a deterministic upper spectral
bound, while the replaced part is Gaussian.  For $R>0$,
$g_R(\theta)^2=2(e^{\theta R}-1-\theta R)/R^2$; the integral definition
also gives $g_0(\theta)=\theta$ and $g_R(0)=0$ without division by zero.
Empty sums cause no change.  In dimension zero both traces vanish.

\paragraph{Reusable declarations and remaining work.}
The snapshot contains sorry-free
\code{NLAlib.gaussian_eq_map_gaussianVector_sqrt_covariance_of_centered},
\code{NLAlib.traceExp_le_traceExp}, and
\code{NLAlib.traceExp_smul_eq_sum}.  They support Gaussian construction
and integrability bounds.  The new finite bookkeeping is the covariance
identity for $Q_A$ and iteration of one-step replacement.  The
BMV consequences \eqref{eq:comparison-bmv-two} remain a genuine prerequisite.

\subsection{Extreme eigenvalues from the Gaussian proxy}
\label{route:gaussian-proxy-eigenvalue-comparison}
\code{gaussian-proxy-eigenvalue-comparison}.
Assume the preceding hypotheses with $d\ge1$, and put
\[
 Y=\sum_iW_i,\quad
 \mu=\mathbb E\lambda_{\max}(Z+\Delta),\quad
 \varphi=\mathbb E\lambda_{\max}(Z),\quad
 s_*^2=\sup_{\norm{u}=1}\operatorname{Var}(u^*Zu),\quad
 V=R\varphi/3+s_*^2.
\]
Here $\varphi\ge0$ by the Rayleigh variational formula and centering.

\paragraph{The exact Lipschitz constant.}
Write $Z=\sum_{j=1}^m g_jA_j$ with independent standard real Gaussian
coordinates and Hermitian $A_j$.  For each unit $u$,
$\sum_j(u^*A_ju)^2=\operatorname{Var}(u^*Zu)\le s_*^2$.
Taking the supremum of the Rayleigh quotients and applying finite
Cauchy--Schwarz proves that
$g\mapsto\lambda_{\max}(\sum_jg_jA_j+\Delta)$ is $s_*$-Lipschitz.
The Gaussian centered MGF inequality therefore gives
\[
 \log\mathbb E e^{a\lambda_{\max}(Z+\Delta)}
 \le a\mu+\frac{a^2s_*^2}{2}\qquad(a\ge0).
\]
Use the MGF form itself, rather than attempting to recover its sharp
constant by integrating a tail inequality.  It is already supplied by
the sorry-free pair
\code{NLAlib.hasEuclideanLogSobolev_gaussianVector} and
\code{NLAlib.hasSubgaussianMGF_of_logSobolev}.

\paragraph{Reduction to a scalar sub-gamma inequality.}
The trace is at most $d$ times the exponential of the largest eigenvalue.
Since $g=g_R(\theta)\ge\theta$, Weyl's inequality gives
\[
 \mathbb E\lambda_{\max}(gZ+\theta\Delta)
 \le\theta\mu+(g-\theta)\varphi.
\]
For $b=R\theta/3<1$, expansion of the exponential and
$(k+2)!\ge2\,3^k$ yield
\[
 g^2\le\frac{\theta^2}{1-b},\qquad
 g-\theta=\frac{g^2-\theta^2}{g+\theta}
 \le\frac{R\theta^2}{6(1-b)}
 \quad(\theta>0).
\]
Combining these estimates with \eqref{eq:comparison-trace-mgf} gives
\begin{equation}\label{eq:comparison-subgamma}
 \log\mathbb E\operatorname{Tr}e^{\theta(Y+\Delta)}
 \le\log d+\theta\mu+
 \frac{V\theta^2}{2(1-c\theta)},\qquad c=R/3.
\end{equation}
Both the expectation and the tail claim now follow from the same scalar
substitution.  For $V>0$ and $s>0$, put
\[
 a=\sqrt{2s/V},\qquad \theta=\frac{a}{1+ca}.
\]
Then
\[
 -\theta(\sqrt{2Vs}+cs)+\frac{V\theta^2}{2(1-c\theta)}=-s.
\]
Exponential Markov applied to the largest eigenvalue proves
\[
 \mathbb P\{\lambda_{\max}(Y+\Delta)
       \ge\mu+\sqrt{2Vs}+cs\}\le de^{-s}.
\]
For the expectation, Jensen and the trace bound give
\[
 \mathbb E\lambda_{\max}(Y+\Delta)
 \le\mu+\frac{\log d}{\theta}
             +\frac{V\theta}{2(1-c\theta)}.
\]
The same substitution with $s=\log d>0$ gives exactly
\[
 \mathbb E\lambda_{\max}(Y+\Delta)
 \le\mu+\sqrt{2V\log d}+\frac R3\log d.
\]
For $d=1$, let $\theta\downarrow0$, or use the scalar centering directly.

\paragraph{Degenerate cases and the tail event.}
If $R=0$, then $W_i\preceq0$ and $\mathbb EW_i=0$ imply $W_i=0$ almost
surely: apply expectation to the nonnegative variable $-\operatorname{Tr}W_i$.
The proxy is zero as well.  The strict event
\[
 \{\lambda_{\max}(Y+\Delta)>
                      \mu+\sqrt{2Vs}+(R/3)s\}
\]
is therefore empty, and its probability is at most $de^{-s}$ for every
$s\ge0$.  A non-strict event at the deterministic mean would be false
for large $s$.  If $V=0$, then $s_*=0$; polarization of finitely many
quadratic forms gives $Z=0$ almost surely.  Matching covariance and
independence then imply all centered $W_i=0$.  For $R>0$ the non-strict
tail event is empty when $s>0$, and the $s=0$ bound is trivial.
Thus the atlas's strict event holds universally; the non-strict event
holds under its stated additional condition $R>0$.
A lower tail follows by applying this theorem to $-W_i,-\Delta$, but
requires the separate lower bound $W_i\succeq-R_-I$.

The sorry-free declarations
\code{NLAlib.integral_lambdaMax_le_log_integral_traceExp_div} and
\code{NLAlib.measureReal_lambdaMax_ge_le_exp_mul_integral_traceExp}
already supply the two Laplace steps.  The only deep unproved dependency
of this route is inherited from the preceding BMV-based replacement.

\subsection{Exchangeable covariance by an antisymmetric difference}
\label{route:exchangeable-matrix-covariance}
\code{exchangeable-matrix-covariance}.
Let $(X,X')$ be exchangeable, square-integrable Hermitian matrices,
$\mu=\mathbb EX$, and assume
$\mathbb E[X-X'\mid X]=\alpha(X-\mu)$ with $\alpha>0$.
For a measurable $F$ with $F(X)$ square-integrable, conditional expectation
and exchangeability give respectively
\[
 \alpha\mathbb E\tr[(X-\mu)F(X)]
       =\mathbb E\tr[(X-X')F(X)],\qquad
 \mathbb E\tr[(X-X')F(X)]
       =-\mathbb E\tr[(X-X')F(X')].
\]
Their average proves
\begin{equation}\label{eq:comparison-exchangeable}
 \mathbb E\tr[(X-\mu)F(X)]
 =\frac1{2\alpha}\mathbb E\tr[(X-X')(F(X)-F(X'))].
\end{equation}
Every product is integrable by finite-dimensional Cauchy--Schwarz.
Real and imaginary parts can be treated separately.  The trace can be
ordinary or normalized, as long as the same convention is used throughout.

Choose real coordinates $x_a$ on the Hermitian space.  Apply the same
argument to $\mathbb E[(x_a-x'_a)(x_b-\mu_b)]$, first using regression
and then swapping.  This proves, for every $a,b$,
\[
 \mathbb E[(x_a-\mu_a)(x_b-\mu_b)]
   =\frac1{2\alpha}\mathbb E[(x_a-x'_a)(x_b-x'_b)].
\]
These finitely many equalities are exactly
\[
 \mathbb E[(X-\mu)\otimes(X-\mu)]
   =\frac1{2\alpha}\mathbb E[(X-X')\otimes(X-X')].
\]
The tensor product is bilinear over the real Hermitian space; it is not
the matrix product $\mathbb E(X-\mu)^2$.

\paragraph{Realizing the pair without conditional laws.}
For $X=\sum_{i=1}^nS_i$, $n\ge1$, with independent square-integrable
$S_i$, take independent copies $S_i'$ and an independent uniform $I$.
Set $X'=X-S_I+S_I'$.  Conditional on $I=i$, swapping $S_i,S_i'$ preserves
the product law; averaging proves exchangeability.  Conditional on all
original summands,
\[
 \mathbb E[X-X'\mid S_1,\ldots,S_n]
   =\frac1n\sum_i(S_i-\mathbb ES_i)=\frac{X-\mathbb EX}{n}.
\]
The tower property yields the required regression with $\alpha=1/n$.
Identical distributions are not required.  A finite product of the
summand laws gives the construction on any underlying probability space.
No analytic matrix-function theorem is needed for this card.
The reference is \href{https://arxiv.org/html/2603.05803v2#S6.SS3}{Tropp,
Section 6.3}; the proof above is the complete coordinate argument.

\subsection{Cauchy-transform universality with constant one}
\label{route:cauchy-transform-gaussian-universality}
\code{cauchy-transform-gaussian-universality}.
Let $d,n\ge1$, let $S_i$ be independent Hermitian matrices with finite
third moments, put $X=\sum_iS_i$, $\mu=\mathbb EX$, and let $Z$ be an
independent Gaussian matrix with the same mean and full covariance.
Fix $\zeta\notin\R$, write $\eta=|\operatorname{Im}\zeta|$, and use
$\tr_d=d^{-1}\operatorname{Tr}$ throughout this subsection.  Define
\[
 R_A=(\zeta I-A)^{-1},\qquad f(A)=R_A^2,\qquad
 M_3=\sum_i\mathbb E\tr_d|S_i-S_i'|^3.
\]
The desired estimate is
\begin{equation}\label{eq:comparison-cauchy-target}
 \left|\mathbb E\tr_d R_X-\mathbb E\tr_d R_Z\right|
 \le\frac{M_3}{\eta^4}.
\end{equation}

\paragraph{Only two finite resolvent identities are needed.}
For Hermitian $A,B,C$ and arbitrary directions $H,K$, set
\begin{align*}
 L(A,B)[H]&=R_A^2HR_B+R_AHR_B^2,\\
 T(A,B,C)[H,K]&=R_A^2HR_BKR_C+R_AHR_B^2KR_C+R_AHR_BKR_C^2.
\end{align*}
Multiplication by the inverses proves $R_A-R_B=R_A(A-B)R_B$.
Use this identity once in each product to obtain
\begin{align}
 f(A)-f(B)&=L(A,B)[A-B],\label{eq:comparison-L}\\
 Df(C)[H]&=L(C,C)[H],\nonumber\\
 L(A,B)[H]-Df(C)[H]
   &=T(A,B,C)[H,B-C]+T(A,C,C)[A-C,H].\label{eq:comparison-T}
\end{align}
For the last equation, insert $L(A,C)[H]$ between the two terms and
telescope first the right argument and then the left argument.
This proves the ordered identity directly.  One must not identify an
ordered second difference with half a mixed Hessian: the latter is the
sum of the two orders, divided by two only after a valid symmetry argument.

\paragraph{A third copy that decouples the covariance.}
Take a second independent copy $S_i''$ and use the same uniform index $I$:
\[
 X'=X-S_I+S_I',\qquad X''=X-S_I+S_I'',\qquad
 D=X-X',\quad E=X'-X'',\quad F=X-X''.
\]
For each fixed $I=i$, $X''$ is independent of $D=S_i-S_i'$ and has
the law of $X$.  As this law does not depend on $i$, $X''$ and $D$ are
also independent without conditioning.  This conclusion is particular
to this resampling construction; an arbitrary exchangeable triple would
not suffice.  The covariance identity from the preceding subsection is
\[
 \operatorname{Cov}_{\otimes}(X)=\frac n2\mathbb E(D\otimes D).
\]

\paragraph{Interpolation and exact cancellation.}
Set
\[
 Y_t=\mu+\sqrt t(X-\mu)+\sqrt{1-t}(Z-\mu),\quad
 u(t)=\mathbb E\tr_d R_{Y_t},\qquad0\le t\le1,
\]
and define $Y_t',Y_t''$ by replacing $X$ with $X',X''$.
Cyclicity of trace gives $D\tr_d R_A[H]=\tr_d(HR_A^2)$, hence
\[
 u'(t)=\frac1{2\sqrt t}\mathbb E\tr_d[(X-\mu)f(Y_t)]
       -\frac1{2\sqrt{1-t}}\mathbb E\tr_d[(Z-\mu)f(Y_t)].
\]
For the first term apply \eqref{eq:comparison-exchangeable} while fixing
$Z$, and then \eqref{eq:comparison-L}; this gives
$\frac n4\mathbb E\tr_d[D L(Y_t,Y_t')[D]]$.
For the second term, ordinary Gaussian integration by parts in real
coordinates gives one half of the covariance paired with $Df(Y_t)$.
By covariance matching, the last covariance identity, and independence
of $D$ from $(X'',Z)$, that term is
$\frac n4\mathbb E\tr_d[D Df(Y_t'')[D]]$.
Thus \eqref{eq:comparison-T} yields the exact six-word identity
\begin{align}\label{eq:comparison-sixwords}
 u'(t)=\frac{n\sqrt t}{4}\mathbb E\tr_d\bigl[
 D T(Y_t,Y_t',Y_t'')[D,E]
 +D T(Y_t,Y_t'',Y_t'')[F,D]\bigr].
\end{align}
This derivation only uses finite coordinate Gaussian integration by parts;
it does not import a general rational-function IBP theorem with unstated
integrability hypotheses.

\paragraph{Each of the six words has the same bound.}
The spectral theorem gives $\norm{R_A}\le\eta^{-1}$ for Hermitian $A$.
For arbitrary matrices, normalized Schatten H\"older and the ideal property
give
\[
 |\tr_d(H_1B_1H_2B_2H_3B_3)|
 \le\prod_{j=1}^3\bigl(\tr_d|H_j|^3\bigr)^{1/3}
                 \prod_{j=1}^3\norm{B_j}.
\]
Ordinary scalar H\"older under expectation now applies even though the
three differences are dependent.  Exchangeability of the three copies
gives
\[
 \mathbb E\tr_d|D|^3=\mathbb E\tr_d|E|^3
   =\mathbb E\tr_d|F|^3=M_3/n.
\]
Every term in \eqref{eq:comparison-sixwords} contains four resolvent
factors and three such differences, so
\[
 |u'(t)|\le\frac{n\sqrt t}{4}\cdot6\cdot
                  \frac{M_3}{n\eta^4}
             =\frac32\sqrt t\,\frac{M_3}{\eta^4}.
\]
Integrate and use $\int_0^1\frac32\sqrt t\,dt=1$ to obtain
\eqref{eq:comparison-cauchy-target} with its exact constant.

\paragraph{Analytic justification and a real-matrix interface.}
On each compact subinterval of $(0,1)$ the derivative of the trace
resolvent is bounded by a constant times $\norm{X-\mu}+\norm{Z-\mu}$,
which is integrable.  The functions $f$ and $Df$ are bounded in norm by
$\eta^{-2}$ and $2\eta^{-3}$ times the direction norm, respectively;
Gaussian IBP applies coordinatewise.  The six-word expression is
integrable by the finite third moments.  At the two endpoints,
$|\tr_d R_{Y_t}|\le\eta^{-1}$ and pointwise continuity imply convergence
by dominated convergence.  Apply the fundamental theorem on
$[a,1-a]$ and let $a\downarrow0$.

The current Schatten product bounds are stated over $\R$.  Their complex
use can be obtained with no loss by the finite realification
\[
 \mathcal R(A)=\begin{pmatrix}\operatorname{Re}A&-\operatorname{Im}A\\
                         \operatorname{Im}A&\operatorname{Re}A\end{pmatrix}.
\]
It preserves products and adjoints; singular values occur twice, so it
preserves the operator norm and normalized Schatten norms.  Also
$\tr_{2d}\mathcal R(A)=\operatorname{Re}\tr_d A$.
Multiply one factor of a trace word by a unit complex scalar making its
trace real and nonnegative, and then realify.  This proves the complex
absolute-trace inequality with exactly the same constant.  The repeated
singular values follow by realifying a complex SVD; the realified unitary
matrices are orthogonal.

\paragraph{Bounded summands.}
If $T_i=S_i-\mathbb ES_i$ satisfies $\norm{T_i}\le L$, then
$\norm{S_i-S_i'}\le2L$ and functional calculus gives
\[
 |S_i-S_i'|^3\preceq2L(S_i-S_i')^2,
 \qquad \mathbb E(S_i-S_i')^2=2\mathbb ET_i^2.
\]
Therefore, with $\sigma^2=\norm{\sum_i\mathbb ET_i^2}$,
\[
 M_3\le4L\tr_d\sum_i\mathbb ET_i^2\le4L\sigma^2.
\]
This proves the advertised $4L\sigma^2/\eta^4$ consequence, including
$L=0$ and $\sigma=0$.

\paragraph{Verified reusable statements.}
All of the following are sorry-free in the inspected snapshot:
\code{NLAlib.integral_mul_eq_sum_covariance_mul_integral},
\code{NLAlib.integral_mul_eq_integral_deriv_gaussianReal},
\code{NLAlib.rationalBlockDifference_sub_left},
\code{NLAlib.rationalBlockDifference_sub_right},
\code{NLAlib.rationalSecondBlockDifference_resolvent},
\code{NLAlib.rationalSecondBlockDifference_mul},
\code{NLAlib.schattenNormReal_mul_le_mul},
\code{NLAlib.schattenNormReal_mul_le_specNorm_mul}, and
\code{NLAlib.schattenNormReal_mul_le_mul_specNorm}.
The first two apply to real and imaginary components separately.  The
explicit formulas above can either bypass the rational-block interface
or be proved by specializing it.  This card has no BMV dependency.
For comparison, see
\href{https://arxiv.org/html/2603.05803v2#S8}{Tropp, Theorem 8.1}.

\subsection{Sample-covariance lower bound from fourth moments}
\label{route:sample-covariance-fourth-moment-lower}
\code{sample-covariance-fourth-moment-lower}.
Let $d\ge2$, let $w\in\R^d$ be centered with finite fourth moments,
and let $K=\mathbb Eww^{\mathsf T}\succ0$.  Assume $\beta\ge1$ and
\[
 \operatorname{Var}(\langle w,a\rangle^2)
       \le\beta^2\bigl(\mathbb E\langle w,a\rangle^2\bigr)^2
       \quad(a\in\R^d).
\]
For independent copies, write
$\widehat K=n^{-1}\sum_iw_iw_i^{\mathsf T}$.
We prove that, for $0<\varepsilon,\delta<1$,
\[
 n\ge\frac{26\beta^2\max\{d,\log(d/\delta)\}}{\varepsilon^2}
 \quad\Longrightarrow\quad
 \mathbb P\{(1-\varepsilon)K\preceq\widehat K\}\ge1-\delta.
\]
This route uses the preceding Gaussian proxy comparison, and therefore
inherits its isolated BMV input.  It does not truncate the observations.

\paragraph{Whitening gives a bounded lower increment for free.}
Put $v=K^{-1/2}w$.  Then $\mathbb Evv^{\mathsf T}=I$, and for unit $u$,
$\operatorname{Var}(\langle v,u\rangle^2)\le\beta^2$.
For
\[
 T_i=\frac{v_iv_i^{\mathsf T}-I}{n},\qquad Y=\sum_iT_i,
\]
we have $\mathbb ET_i=0$ and $T_i\succeq-I/n$, although $\norm{v_i}$
need not be bounded.  The fourth moments guarantee finite second moments
of $T_i$.  Its centered Gaussian proxy $Z$ satisfies
\[
 \operatorname{Var}(u^{\mathsf T}Zu)
 =\frac1n\operatorname{Var}(\langle v,u\rangle^2)
 \le\beta^2/n\qquad(\norm{u}=1).
\]
Taking the supremum gives $s_*^2\le\beta^2/n$.

\paragraph{A finite net estimate with the required numerical constant.}
Take a $1/8$-net $\mathcal N$ on the real unit sphere of size at most
$17^d$.  For symmetric $Z$,
\[
 \norm{Z}\le\frac43\max_{u\in\mathcal N}|u^{\mathsf T}Zu|.
\]
Indeed, approximate a unit vector maximizing the absolute quadratic
form; the discrepancy of the two forms is at most
$2(1/8)\norm{Z}$, which can be moved to the left side.
Each scalar form is centered Gaussian with variance at most $\beta^2/n$.
For any finite collection of centered Gaussians of variance at most
$v^2$, the bound
$\mathbb E\max_j|g_j|\le v\sqrt{2\log(2N)}$ follows from
$\mathbb Ee^{t\max_j|g_j|}\le2N e^{t^2v^2/2}$ and optimization;
no independence of the forms is used.  Thus
\[
 \mathbb E\norm Z\le\frac43\frac\beta{\sqrt n}
       \sqrt{2\log(2\cdot17^d)}
       \le a:=\sqrt{12\beta^2d/n}.
\]
For the last inequality, $d\ge2$, $\log17<3$, and $\log2<3/4$ imply
$\log(2\cdot17^d)\le27d/8$.  The two logarithmic bounds are certified
by positive finite partial sums for $e^3>17$ and $e^{3/4}>2$.

\paragraph{The failure event and arithmetic certificate.}
Apply the upper proxy theorem to $-T_i$, with $R=1/n$, or equivalently
its lower-tail version to $T_i$.  Gaussian symmetry gives
$\mathbb E\lambda_{\min}(Z)=-\mathbb E\lambda_{\max}(Z)\ge-a$.
For $s=\log(d/\delta)>0$, the probability of
\[
 \lambda_{\min}(Y)
 \le-\left[a+\sqrt{\frac{2s}{n}(\beta^2+a/3)}+\frac{s}{3n}\right]
\]
is at most $de^{-s}=\delta$.  This is the lower-tail failure event;
reversing its inequality would not express the desired conclusion.
Under the sample-size assumption, the three error terms satisfy
\[
 a\le\frac{53}{78}\varepsilon,\qquad
 \sqrt{\frac{2s}{n}(\beta^2+a/3)}\le\frac{24}{78}\varepsilon,
 \qquad \frac{s}{3n}\le\frac1{78}\varepsilon.
\]
For a purely rational check, the first inequality uses
$a^2\le(6/13)\varepsilon^2$ and $53^2\ge2808$.
For the second, use $\varepsilon\le1$, $\beta\ge1$, and the first bound:
\[
 \frac{2s}{n}(\beta^2+a/3)
 \le\frac{\varepsilon^2}{13}
          \left(1+\frac{53}{234}\right)
 =\frac{574}{78^2}\varepsilon^2
 \le\frac{24^2}{78^2}\varepsilon^2.
\]
The third follows from $s/n\le\varepsilon^2/(26\beta^2)$.
The terms sum to at most $\varepsilon$.  Hence
$\mathbb P\{\lambda_{\min}(Y)<-\varepsilon\}\le\delta$;
on its complement $I+Y\succeq(1-\varepsilon)I$.
Congruence by $K^{1/2}$ gives the required event for $\widehat K$.
Positive definiteness of $K$ is exactly what justifies whitening.

The net interfaces are already sorry-free:
\code{NLAlib.exists_finset_sphere_net} and
\code{NLAlib.specNorm_le_div_of_net_quadForm}.
The reference target is
\href{https://arxiv.org/html/2603.04365v5#S3.SS3}{Tropp, Theorem 3.3}.
The finite net and the displayed rational inequalities account for the
constant $26$ explicitly; the result is one-sided covariance preservation,
not a two-sided error guarantee.

\section*{Implementation order and audit record}
\addcontentsline{toc}{section}{Implementation order and audit record}
The highest-value shared new lemma is the finite Hermite moment bound:
it feeds the partial-transpose card, covariance-sensitive norm card, and
both chaos norm estimates. Its proof uses only finitely many Gaussian
polynomial moments, sparse finite matrices, and the coefficient-block
order bound. The scalar domination lemma and the coefficient-reindexing
identities should be packaged separately so that all chaos proofs reuse them.

The next small interfaces are support functions plus compact nets, the
exchangeable-pair coordinate identity, realification of complex Schatten
norms, and the six ordered resolvent words. They provide the width and
universality results without importing a larger process, rational-function,
or complex integration framework. The Fourier card needs finite Schur
averaging and an explicit orthonormal coefficient basis; it does not need
infinite-dimensional harmonic analysis.

The remaining deep inputs should remain explicit dependencies. A
covariance-sensitive norm bound does not imply two-sided spectral-set
comparison, local invertibility of a matrix transform does not identify
its physical global branch, and a positive second trace-exponential
derivative does not establish the higher positivity supplied by BMV.
The document keeps these distinctions throughout.

Independent mathematical checks were performed across the finite Hermite,
chaos, Gaussian width, Fourier, exchangeable-pair, resolvent, and variational
reductions. The checks include the Gaussian variance conventions, normalized
versus unnormalized traces, the Cauchy-transform constant $1$, the sample-size
constant $26$, and all stated degenerate cases. No numerical experiment is
used as proof. The exact original contracts and their primary source labels
remain available in the
\href{https://github.com/tchen-research/NLAlib/blob/faef0a270d762f0bf0e90603544d50f289873d0e/atlas/atlas.json}
{dated atlas snapshot}; source links are also placed beside the derivations.
\end{document}
